\documentclass[12pt,a4paper]{article}

% Encoding and fonts
\usepackage[utf8]{inputenc}
\usepackage[T1]{fontenc}
\usepackage{lmodern}

% Page layout
\usepackage[top=1in, bottom=1in, left=1in, right=1in]{geometry}

% Typography
\usepackage{microtype}
\usepackage{parskip}

% Language
\usepackage[english]{babel}

% Headers and footers
\usepackage{fancyhdr}
\pagestyle{fancy}
\fancyhf{}
\fancyhead[L]{\leftmark}
\fancyhead[R]{\thepage}
\fancyfoot[C]{\thepage}
\renewcommand{\headrulewidth}{0.4pt}
\renewcommand{\footrulewidth}{0pt}

% Section styling
\usepackage{titlesec}
\titleformat{\section}{\Large\bfseries}{\thesection}{1em}{}
\titleformat{\subsection}{\large\bfseries}{\thesubsection}{1em}{}
\titleformat{\subsubsection}{\normalsize\bfseries}{\thesubsubsection}{1em}{}
\titlespacing*{\section}{0pt}{2.5ex plus 1ex minus .2ex}{1.5ex plus .2ex}
\titlespacing*{\subsection}{0pt}{2ex plus 1ex minus .2ex}{1ex plus .2ex}
\titlespacing*{\subsubsection}{0pt}{1.5ex plus 1ex minus .2ex}{0.8ex plus .2ex}

% List formatting
\usepackage{enumitem}
\setlist[itemize]{topsep=0.5ex, itemsep=0.25ex, parsep=0pt}
\setlist[enumerate]{topsep=0.5ex, itemsep=0.25ex, parsep=0pt}

% Essential packages
\usepackage{graphicx}
\usepackage[table]{xcolor}
\usepackage{booktabs}
\usepackage{array}
\usepackage{tabularx}
\usepackage{longtable}
\usepackage{amsmath}
\usepackage{amssymb}
\usepackage[normalem]{ulem}
\rowcolors{2}{gray!10}{white}
\renewcommand{\arraystretch}{1.3}

% Cross-references
\usepackage{hyperref}
\usepackage{cleveref}

% Hyperref setup
\hypersetup{
    colorlinks=true,
    linkcolor=blue,
    filecolor=magenta,
    urlcolor=cyan,
}

% Code listings
\usepackage{listings}
\definecolor{codebg}{RGB}{248,248,248}
\definecolor{codeframe}{RGB}{200,200,200}
\definecolor{codekeyword}{RGB}{0,0,180}
\definecolor{codecomment}{RGB}{0,128,0}
\definecolor{codestring}{RGB}{163,21,21}
\lstset{
    basicstyle=\ttfamily\small,
    breaklines=true,
    breakatwhitespace=false,
    frame=single,
    framerule=0.5pt,
    rulecolor=\color{codeframe},
    backgroundcolor=\color{codebg},
    keepspaces=true,
    showstringspaces=false,
    tabsize=4,
    xleftmargin=1em,
    xrightmargin=1em,
    aboveskip=1em,
    belowskip=1em,
    keywordstyle=\color{codekeyword}\bfseries,
    commentstyle=\color{codecomment}\itshape,
    stringstyle=\color{codestring},
}

% YAML language definition for listings
\lstdefinelanguage{YAML}{
    keywords={true,false,null,y,n},
    keywordstyle=\color{blue}\bfseries,
    sensitive=false,
    comment=[l]{\#},
    morecomment=[s]{/*}{*/},
    commentstyle=\color{gray}\ttfamily,
    stringstyle=\color{red}\ttfamily,
    morestring=[b]',
    morestring=[b]",
}

% TOML language definition for listings
\lstdefinelanguage{TOML}{
    keywords={true,false},
    keywordstyle=\color{blue}\bfseries,
    sensitive=true,
    comment=[l]{\#},
    commentstyle=\color{gray}\ttfamily,
    stringstyle=\color{red}\ttfamily,
    morestring=[b]',
    morestring=[b]",
}

% Dockerfile language definition for listings
\lstdefinelanguage{Dockerfile}{
    keywords={FROM,RUN,CMD,LABEL,EXPOSE,ENV,ADD,COPY,ENTRYPOINT,VOLUME,USER,WORKDIR,ARG,ONBUILD,STOPSIGNAL,HEALTHCHECK,SHELL},
    keywordstyle=\color{blue}\bfseries,
    sensitive=false,
    comment=[l]{\#},
    commentstyle=\color{gray}\ttfamily,
    stringstyle=\color{red}\ttfamily,
    morestring=[b]',
    morestring=[b]",
}

% JSON language definition for listings
\lstdefinelanguage{JSON}{
    keywords={true,false,null},
    keywordstyle=\color{blue}\bfseries,
    sensitive=true,
    commentstyle=\color{gray}\ttfamily,
    stringstyle=\color{red}\ttfamily,
    morestring=[b]",
}

% Code listing float environment
\usepackage{float}
\newfloat{listing}{htbp}{lol}[section]
\floatname{listing}{Listing}

% Admonition boxes
\usepackage{tcolorbox}
\tcbuselibrary{skins,breakable}

% Admonition style definitions
\newtcolorbox{admonitionbox}[2][]{
    enhanced,
    breakable,
    colback=#2!5,
    colframe=#2!80!black,
    fonttitle=\bfseries,
    title=#1,
    left=2mm,
    right=2mm,
}

% Synthon tuple boxes
\newtcolorbox{synthonbox}{
    enhanced,
    colback=blue!5,
    colframe=blue!75!black,
    fontupper=\large,
    center upper,
    boxrule=1pt,
    arc=4pt,
}

\begin{document}

\title{SynthOmnicon --- Primitive-Derived Predictions}
\date{\today}

\maketitle

\emph{Framework version: v0.4.3 $\cdot$ Date: 2026-03-17}

This document is a living ledger of every prediction the SynthOmnicon framework has generated \textbf{purely from primitive assignments} --- with no domain-specific physics inserted --- and the current experimental or computational status of each. Three tiers are distinguished:

\begin{itemize}
    \item \textbf{Tier I --- Experimentally Confirmed:} The prediction follows from ordinal/structural primitives alone; experimental measurement confirms it.
    \item \textbf{Tier II --- Computationally Validated:} The prediction is confirmed against high-quality DFT/SAPT benchmarks or literature values; no independent wet-lab measurement is claimed.
    \item \textbf{Tier III --- Falsifiable, Pending Test:} A specific, measurable quantity is predicted from primitives; experimental data does not yet exist or has not been found.
\end{itemize}

The eleven-primitive tuple is \texttt{\langleD; T; R; P; F; K; G; \Gamma; Phi; S; \Omega\rangle}. "Purely primitive-derived" means the prediction follows from the ordinal lattice, the composition axioms, or the algebraic operations (tensor, meet, join, lift, path) applied to the encoded primitives --- with no domain equations substituted.

\begin{center}
\rule{0.5\textwidth}{0.4pt}
\end{center}

\subsection{Tier I --- Experimentally Confirmed}

\subsubsection{P-1 $\cdot$ F-floor ratchet (CB{[}7{]} competitive displacement)}

\textbf{Primitive basis.} F is a totally ordered primitive: $F_{\hbar}$ \textgreater{} F\_eth \textgreater{} $F_{\ell}$. The HotSwap ratchet (Axiom-adjacent constraint) states that a guest can displace a bound guest if and only if its F tier is strictly higher (F-floor hard constraint). From F ordinals alone, six directional predictions follow:

\begin{table}[htbp]
\centering
\small
\rowcolors{1}{white}{white}
\begin{tabularx}{\textwidth}{>{\centering\arraybackslash}X >{\centering\arraybackslash}X >{\centering\arraybackslash}X}
\toprule
\rowcolor{gray!25}\textbf{Displacement event} & \textbf{F prediction} & \textbf{Outcome} \\
\midrule
\rowcolors{1}{white}{gray!10}
Fc displaces Ad & $F_{\hbar}$ \textgreater{} F\_eth $\rightarrow$ allowed &  confirmed \\
Fc displaces DABCO & $F_{\hbar}$ \textgreater{} $F_{\ell}$ $\rightarrow$ allowed &  confirmed \\
Ad displaces DABCO & F\_eth \textgreater{} $F_{\ell}$ $\rightarrow$ allowed &  confirmed \\
Ad displaces Fc & F\_eth \textless{} $F_{\hbar}$ $\rightarrow$ blocked &  confirmed \\
DABCO displaces Ad & $F_{\ell}$ \textless{} F\_eth $\rightarrow$ blocked &  confirmed \\
DABCO displaces Fc & $F_{\ell}$ \textless{} $F_{\hbar}$ $\rightarrow$ blocked &  confirmed \\
\bottomrule
\end{tabularx}
\end{table}

\textbf{A seventh, meta-level prediction:} the ratchet is strictly asymmetric --- displacement is irreversible at each F boundary. No higher-F guest can be ejected by a lower-F competitor under any concentration condition. This is a topological claim about the F lattice: there is no path from a higher-F occupied state to a lower-F occupied state by competitive displacement alone. The experimental confirmation is that \emph{none} of the six systems shows reversal even under large excess of the weaker guest (Ad and DABCO were tested at excess; Fc was not displaced).

\textbf{Experimental anchor.} Kim JACS 2001; Assaf \& Nau \emph{Chem. Soc. Rev.} 2015; Sindelar JOC 2007.
CB{[}7{]}$\cdot$Fc: K\_a = 3 $\times$ 10¹$^{2}$ M⁻¹ ($F_{\hbar}$). CB{[}7{]}$\cdot$Ad: K\_a = 4 $\times$ 10⁸ M⁻¹ (F\_eth). CB{[}7{]}$\cdot$DABCO: K\_a = 2 $\times$ 10⁵ M⁻¹ ($F_{\ell}$). \emph{$F_{\ell}$ tier activated March 17, 2026 --- threshold K\_a \textless{} 10⁷ M⁻¹ formalised from the DABCO binding data at this date; all six directional predictions confirmed after tier activation.}

\textbf{Score: 6/6 directional predictions + asymmetric ratchet topology confirmed from ordinal F ranking alone.}

Validation row: V-1 $\cdot$ \xi\_CP: not applicable (qualitative ordering).

\begin{center}
\rule{0.5\textwidth}{0.4pt}
\end{center}

\subsubsection{P-2 $\cdot$ Soai autocatalytic amplification $\rightarrow$ Frank bifurcation (Factor 7)}

\textbf{Primitive basis.} The Frank-model criticality fingerprint (Factor 7) fires when the four co-requisites D\_$\infty$ + T\_⋈ + P\_directional + $F_{\hbar}$ are simultaneously present, yielding $\phi$\_c score \textgreater{} 0.3 and candidacy score 0.920. This is derived purely from the primitive co-occurrence rule --- no kinetic or mechanism-specific information is inserted.

\textbf{Prediction.} A synthon satisfying D\_$\infty$ (temporal/autocatalytic) + T\_⋈ (bowtie/cyclic closure) + P\_DA (donor-acceptor directionality) + $F_{\hbar}$ (high fidelity) will exhibit spontaneous symmetry breaking at ee = 0, following the pitchfork bifurcation topology of the Frank model.

\textbf{Experimental anchor.} Soai reaction (pyrimidyl alkanol, Zn-mediated): Soai \emph{JACS} 1995; Gridnev \emph{Angew. Chem.} 2010; Shibata \emph{JACS} 2009.
Active species: {[}Zn₂$\cdot$(pyrimidylalkoxide)₂$\cdot$iPr₂Zn{]} dimer. $\Delta$G‡ = 14.9 kcal/mol (62.3 kJ/mol); $F_{\hbar}$, K\_mod.
Varma ratio \xi\_r/ln(\x$i_{\tau}$) = 0.94 $\approx$ 1.0 --- approaching $\Phi$\_c. Factor 7 correctly identifies this as the highest-confidence $\Phi$\_c candidate in the catalog. \emph{Note: Factor 7 firing and Tier I status apply after continuous-reset grounding was added to the Soai catalog entry (Axiom 6 compliance, March 17, 2026). Prior to grounding, the entry lacked a named driving gradient and was flagged by Pass 1 audit. The Factor 7 co-requisites (D\_$\infty$ + T\_⋈ + P\_DA + $F_{\hbar}$) are only valid once D\_$\infty$ is grounded.}

Validation row: V-3 $\cdot$ Candidacy score: 0.920 $\cdot$ \xi\_CP: \textasciitilde{}9.21 nats (estimated D\_$\infty$ tier).

\begin{center}
\rule{0.5\textwidth}{0.4pt}
\end{center}

\subsubsection{P-3 $\cdot$ Proline-aldol ee prediction from F\_cycle}

\textbf{Primitive basis.} The proline-catalyzed aldol cycle is assigned D\_$\infty$ (temporal), F\_eth (moderate fidelity), K\_mod. From F\_eth and the operative $\Delta$G‡, the Eyring-based fidelity per turn is F\_cycle $\approx$ 0.999--0.9999. The facial selectivity (re/si discrimination) follows from the Zimmermann--Traxler chair TS geometry encoded at the F\_eth tier: $\Delta$$\Delta$G‡\_si-re = 5--8 kJ/mol, yielding ee = \textbf{70--85\%} from primitives alone.

\textbf{Experimental anchor.} Blackmond RPKA 2004; Houk/List DFT M06-2X/6-31+G(d,p) (2004). Acetone + 4-nitrobenzaldehyde, DMSO: \textbf{74\% ee} measured.

\textbf{The framework's first quantitatively grounded cross-domain prediction tied to a measured stereochemical outcome.} Predicted range (70--85\%) brackets the experimental value (74\%) without any fitting.

Also confirmed: Varma probe ratio = 0.189 ≪ 1.0 $\rightarrow$ $\Phi$\_sub as predicted (not critical). Structural interpretation: spatial correlation length $\geq$ 60 Å would be required for criticality --- inconsistent with sub-molecular enamine geometry.

Live calibration (\texttt{syncon info-bits --calibrate}): I\_rec = 7.98 bits $\cdot$ I\_net = 6.61 bits $\cdot$ I+solvent = 12.57 bits (within expected range 6.0--9.0 bits for I\_rec). 

Validation row: V-2 $\cdot$ \xi\_CP: 9.21 {[}9.09--9.36{]} nats.

\begin{center}
\rule{0.5\textwidth}{0.4pt}
\end{center}

\subsubsection{P-4 $\cdot$ Ice VI multi-ordering-landscape prediction (K\_fast causal encoding)}

\textbf{Primitive basis.} Ice VI is assigned K\_fast ($\Delta$G‡ \textless{} 60 kJ/mol for proton reorientation). From K\_fast alone, the framework predicts that the system can explore multiple ordering landscapes during cooling --- more than one distinct ordered descendant phase should exist. Additionally, a K\_fast $\rightarrow$ K\_slow flip (a single primitive change) is predicted to coincide with the ordering phase transition.

\textbf{Experimental anchor.} Yamane \emph{et al.} (2021) and Gasser \emph{et al.} (2021) discover ice XIX (\textgreater{}1.5 GPa) as a second ordered descendant from ice VI, complementing ice XV (\textasciitilde{}1.0 GPa, Salzmann 2009). Dielectric relaxation measurements (Yamane 2021) directly confirm the K\_fast assignment: ice VI relaxation time is orders of magnitude shorter than those of ice XV or XIX.

\textbf{Three predictions confirmed from a single primitive assignment:}

\begin{enumerate}
    \item Multiple ordered descendants of ice VI exist --- confirmed (ice XV and ice XIX).
    \item K\_fast $\rightarrow$ K\_slow is the ordering transition --- confirmed (the ordered phases encode K\_slow).
    \item Deep-glassy states accessible under rapid cooling --- consistent with Rosu-Finsen \& Salzmann 2020 partially ordered states.
\end{enumerate}

The framework encoded K\_fast \emph{without} any access to the dielectric data.

Validation row: V-4 $\cdot$ \xi\_CP: qualitative (extended network phase; $\Delta$G definition non-trivial).

\begin{center}
\rule{0.5\textwidth}{0.4pt}
\end{center}

\subsubsection{P-5 $\cdot$ Fidelity ordering: carboxylic acid \textgreater{} amide (H-bond dimer series, Transformation \#1)}

\textbf{Primitive basis.} Cyclic R₂$^{2}$(8) dimers are assigned P\_$\pm$ and F ordered by the strength of the H-bond donor: carboxylic acid (AA$\cdot$AA, P\_$\pm$\^{}\psi, $F_{\hbar}$) \textgreater{} mixed heterodimer (AA$\cdot$amide, F\_eth) \textgreater{} amide homodimer (F\_eth). This ordering follows from the primitive assignment conventions alone --- no $\Delta$E values were input.

\textbf{Computational confirmation.} B3LYP-D3(BJ)/6-311+G(d,p) + BSSE:
$\Delta$E: --64.2 kJ/mol (AA$\cdot$AA) \textgreater{} --51.8 kJ/mol (AA$\cdot$amide) \textgreater{} --39.6 kJ/mol (amide$\cdot$amide).
$\Delta$G₂₉₈ (gas): \textasciitilde{}--12 / \textasciitilde{}--10 / \textasciitilde{}--8 kJ/mol. Fidelity ratio \textasciitilde{}1.9. CSD propensity data confirms the ordering. Live calibration for AA$\cdot$AA homodimer (\texttt{syncon info-bits --calibrate}): I\_rec = 9.39 bits $\cdot$ I\_net = 8.02 bits $\cdot$ I+solvent = 13.98 bits (within expected range 9.0--10.5 bits for I\_rec). 

\xi\_CP: 6.66 {[}6.56--6.77{]} / 8.19 {[}8.07--8.32{]} / 8.70 {[}8.56--8.86{]} nats (descending F $\rightarrow$ ascending cost).

Validation row: \#1.

\begin{center}
\rule{0.5\textwidth}{0.4pt}
\end{center}

\subsubsection{P-6 $\cdot$ Triple H-bond induction superlinearity (G\_beth $\rightarrow$ G\_gimel, Transformation \#5)}

\textbf{Primitive basis.} Axiom 3 (cooperative amplification) predicts that when a network of contacts crosses the cooperative threshold, the induction component of binding becomes superlinear in the number of contacts. The G\_beth $\rightarrow$ G\_gimel transition (local $\rightarrow$ mesoscale granularity) is the primitive encoding of this crossing. The prediction: at the triple H-bond array, induction should contribute 2.5--3.5$\times$ its single-contact value, and $\Delta$E should not be strictly additive.

\textbf{Computational confirmation.} SAPT2+/aug-cc-pVDZ on single/double/triple H-bond arrays:
$\Delta$E: \textasciitilde{}--30 / \textasciitilde{}--60 / \textasciitilde{}--95--110 kJ/mol. Induction fraction: \textasciitilde{}10--15\% (single) $\rightarrow$ \textasciitilde{}30--40\% (triple), 2.5--3.5$\times$ superlinear increase confirmed. Cooperativity factor 1.25 (literature range 1.2--1.4). Live calibration (\texttt{syncon info-bits --calibrate}): I\_rec = 16.57 bits $\cdot$ I\_net = 15.19 bits $\cdot$ I+solvent = 21.15 bits (within expected range 14.0--18.0 bits for I\_rec). 

The superlinearity is the direct computational signature of the G\_beth$\rightarrow$G\_gimel primitive transition encoded before the computation was run.

\xi\_CP: 7.65 {[}7.59--7.72{]} nats (triple array). Validation row: \#5.

\begin{center}
\rule{0.5\textwidth}{0.4pt}
\end{center}

\subsubsection{P-7 $\cdot$ Fidelity ordering: halogen bond \textgreater{} chalcogen bond (Transformation \#7)}

\textbf{Primitive basis.} 4-Iodopyridine I$\cdot$$\cdot$$\cdot$N halogen bond is assigned F\_eth; 4-(methylthio)pyridine S$\cdot$$\cdot$$\cdot$N chalcogen bond is assigned $F_{\ell}$. The F ordering F\_eth \textgreater{} $F_{\ell}$ predicts: (i) higher $\Delta$E for halogen; (ii) larger ESP $\sigma$-hole depth; (iii) narrower angular window (SELECTIVE grammar) for halogen vs. broader (BROAD grammar) for chalcogen.

\textbf{Computational confirmation.} B3LYP-D3(BJ)/6-311+G(d,p) + BSSE; Multiwfn ESP:
$\Delta$E: --28.4 kJ/mol (I$\cdot$$\cdot$$\cdot$N) vs. --14.9 kJ/mol (S$\cdot$$\cdot$$\cdot$N). Fidelity ratio: \textbf{1.91}. V\_max: +165 kJ/mol (iodine) vs. +105 kJ/mol (sulfur). Angular window: C--I$\cdot$$\cdot$$\cdot$N $\pm$2.5$^{\circ}$ (halogen) vs. \textasciitilde{}$\pm$12$^{\circ}$ (H-bond) --- confirming SELECTIVE grammar.

Both the energetic ordering and the grammar assignment follow from the primitive encoding before any DFT is run.

\xi\_CP: 7.59 {[}7.47--7.73{]} nats (halogen dimer) $\cdot$ 8.40 {[}8.31--8.49{]} nats (trimer). Validation row: \#7.

\begin{center}
\rule{0.5\textwidth}{0.4pt}
\end{center}

\subsection{Tier II --- Computationally Validated}

\subsubsection{P-8 $\cdot$ Chelate granularity amplification (G\_beth $\rightarrow$ G\_aleph, Transformation \#3)}

\textbf{Primitive basis.} Switching from two monodentate pyridines (two separate binding events, G\_beth) to one bidentate bipyridine (single event that constrains the whole coordination sphere, G\_aleph) is a G-primitive upgrade. The framework predicts a step increase in cumulative binding energy and thermodynamic chelate gain, driven by preorganisation at G\_aleph.

\textbf{Computational confirmation.} B3LYP-D3(BJ)/6-311+G(d,p): {[}Zn(py)₂Cl₂{]} --263.1 kJ/mol vs. {[}Zn(bpy)Cl₂{]} --312.6 kJ/mol. Chelate gain: \textasciitilde{}49 kJ/mol gas; \textasciitilde{}60--90 kJ/mol solvated. Consistent with experimental K\_chelate/K\_mono$^{2}$ \textasciitilde{} 10$^{2}$--10⁴.

\xi\_CP: \textasciitilde{}9.0 $\rightarrow$ \textasciitilde{}7.5 nats (G\_beth $\rightarrow$ G\_aleph raises constraint propagation efficiency).

\begin{center}
\rule{0.5\textwidth}{0.4pt}
\end{center}

\subsubsection{P-9 $\cdot$ Dynamic covalent D\_$\infty$ emergence (imine, Transformation \#4)}

\textbf{Primitive basis.} A static covalent bond ($R_{\subseteq}$, $F_{\hbar}$, K\_slow) becomes a dynamic covalent bond ($R_{\subseteq}$+‡, F\_eth, K\_mod) when the barrier drops below \textasciitilde{}60 kJ/mol in aqueous solution, enabling error-correction through hydrolysis/re-condensation. This is a D\_$\infty$ character assignment from a K primitive threshold alone.

\textbf{Computational confirmation.} B3LYP-D3/6-311+G(d,p): forward imine condensation endergonic +38.7 kJ/mol (gas); barrier \textasciitilde{}162 kJ/mol (gas), \textasciitilde{}90--120 kJ/mol (aq.). Confirms K\_slow $\rightarrow$ K\_mod crossing at the solvation boundary. Same F\_eth tier, distinct K --- the framework correctly distinguishes two operationally different systems under a shared thermodynamic classification.

\begin{center}
\rule{0.5\textwidth}{0.4pt}
\end{center}

\subsubsection{P-10 $\cdot$ Axiom 1 as classical boundary detector (quantum particle series, V-5)}

\textbf{Primitive basis.} Axiom 1 states T\_⋈ + P\_$\pm$ $\rightarrow$ F $\geq$ F\_eth (cyclic self-complementary motifs amplify fidelity above $F_{\ell}$ in the classical domain). For entangled spin, the LLM correctly assigned T\_⋈ + P\_$\pm$\^{}sym but erroneously assigned $F_{\ell}$ --- a violation that persisted through 3 refinement iterations.

\textbf{Framework prediction.} The violation is irresolvable in the classical frame because F (constraint reliability) and Shannon capacity decouple in the quantum domain. The pattern T\_⋈ + P\_$\pm$ + $F_{\ell}$ is a quantum boundary signature: classically unreachable, quantum-mechanically reachable.

\textbf{Confirmed numerically.} Corrected spin tuple: \langleD\_$\land$; T\_⋈; $R_{\supseteq}$; P\_$\pm$\^{}sym; $F_{\hbar}$; K\_trap; G\_ℵ; $\Gamma$\_$\land$(SELECTIVE); $\Phi$\_sub\rangle --- Axiom 1 satisfied. Spin singlet fires with 100\% reliability $\rightarrow$ $F_{\hbar}$, not $F_{\ell}$. The Axiom 1 violation is a domain diagnostic, not a falsification of the axiom.

Five quantum systems encoded from primitive definitions alone produced five distinct, physically defensible tuples with no chemical template available (V-5).

\begin{center}
\rule{0.5\textwidth}{0.4pt}
\end{center}

\subsubsection{P-11 $\cdot$ $\Omega$ lattice semantics (tensor, meet, join on topological protection classes)}

\textbf{Primitive basis.} $\Omega$ is an ordered lattice: TRIVIAL \textless{} ℤ₂ \textless{} ℤ \textless{} NON-ABELIAN. Tensor inherits the stronger protection (max); meet takes the conservative guarantee (min); join gives the capability ceiling (max). These rules follow from the lattice definition alone.

\textbf{Confirmed numerically} (v0.4.0 catalog, $\Omega$ operations):

\begin{table}[htbp]
\centering
\small
\rowcolors{1}{white}{white}
\begin{tabularx}{\textwidth}{>{\centering\arraybackslash}X >{\centering\arraybackslash}X >{\centering\arraybackslash}X}
\toprule
\rowcolor{gray!25}\textbf{Operation} & \textbf{Result} & \textbf{Physical interpretation} \\
\midrule
\rowcolors{1}{white}{gray!10}
spin ⊓ kitaev\_chain & $\Omega$\_TRIVIAL & Singlet has no topological protection \\
fqh ⊔ TI & $\Omega$\_NON\_ABELIAN & Non-Abelian dominates capability ceiling \\
tensor(kitaev\_chain, qubit) & $\Omega$\_Z + $F_{\ell}$ & Chain protects topology; qubit interface degrades fidelity \\
tensor(fqh, kitaev\_chain) & $\Omega$\_NON\_ABELIAN & Non-Abelian invariant dominates \\
\bottomrule
\end{tabularx}
\end{table}

\texttt{tensor(kitaev\_chain, qubit)} specifically predicts the practical challenge of topological qubits: $\Omega$\_Z (topological protection) but $F_{\ell}$ (fidelity bottleneck at the qubit interface). The framework identifies F --- not $\Omega$ --- as the load-bearing primitive for qubit quality.

\begin{center}
\rule{0.5\textwidth}{0.4pt}
\end{center}

\subsubsection{P-12 $\cdot$ K\_trap $\rightarrow$ K\_MBL universality (+2.303 nat across all gap-protected topological phases)}

\textbf{Primitive basis.} K\_trap encodes a coherent many-body gap. K\_MBL encodes a disorder-frozen many-body localized state. The K ordinal structure assigns a universal \xi\_CP cost to this transition. Perturbation sweeps across all three gap-protected quantum synthons (spin\_singlet, kitaev\_chain, fqh\_moore\_read) yield:

\begin{table}[htbp]
\centering
\small
\rowcolors{1}{white}{white}
\begin{tabularx}{\textwidth}{>{\centering\arraybackslash}X >{\centering\arraybackslash}X >{\centering\arraybackslash}X}
\toprule
\rowcolor{gray!25}\textbf{Synthon} & \textbf{K shift} & \textbf{$\Delta$\xi\_CP} \\
\midrule
\rowcolors{1}{white}{gray!10}
spin\_singlet & TRAP $\rightarrow$ MBL & +2.303 nats \\
kitaev\_chain & TRAP $\rightarrow$ MBL & +2.303 nats \\
fqh\_moore\_read & TRAP $\rightarrow$ MBL & +2.303 nats \\
topological\_insulator & SLOW $\rightarrow$ MOD & $-$0.847 nats \\
\bottomrule
\end{tabularx}
\end{table}

The +2.303 nat cost is identical across all three regardless of their different $\Omega$ classes (ℤ, ℤ\_class via spin, non-Abelian). The TI shows the opposite sign because K\_slow $\rightarrow$ K\_mod is decreasing the kinetic barrier, reducing cost.

\textbf{This is a purely ordinal-arithmetic result} --- no disorder physics was input. The universality emerges from K's ordinal structure alone.

See Tier III P-15 for the experimental falsifiability statement.

\begin{center}
\rule{0.5\textwidth}{0.4pt}
\end{center}

\subsubsection{P-13 $\cdot$ Phase boundary at d $\approx$ 9.52 (topological matter / quantum particle split)}

\textbf{Primitive basis.} The SynthOmnicon tuple distance (weighted Euclidean over all primitives) applied to the 8-member quantum catalog, followed by Ward hierarchical clustering, produces a two-branch dendrogram with primary separation at d $\approx$ 9.52 --- without any physics being input to the distance metric.

\textbf{Interpretation confirmed numerically.} Branch 1 (extended topological matter: fqh, TI) vs. Branch 2 (quantum particles + engineered systems). The split tracks the D primitive: D\_△ systems (collective 2D/3D ground states) vs. D\_$\land$/D\_$\infty$ systems (point particles and chains). Proton $\leftrightarrow$ electron: d = 1.80 (only P differs) --- the framework's closest pair prediction, confirmed by atomic physics.

The phase diagram is generated from syntax alone via \texttt{syncon phase-diagram}. No physics is inserted between the primitive encoding and the dendrogram.

\begin{center}
\rule{0.5\textwidth}{0.4pt}
\end{center}

\subsubsection{P-14 $\cdot$ Rotaxane steric cliff: K\_mod vs. K\_trap from sub-Å stopper geometry (Transformation \#8)}

\textbf{Primitive basis.} The mechanical bond R\_⇌ encodes a topological control mechanism: the steric cliff (sub-Å methyl repositioning flips the barrier \textgreater{}5$\times$) is a direct consequence of the aperture constraint, not a continuous Morse potential. K\_mod (slippage-enabled) vs. K\_trap (locked) are predicted to differ by a discontinuous barrier at constant R\_⇌ and T\_⋈.

\textbf{Partially anchored via literature proxy.} Groppi \emph{et al.} (\emph{Angew. Chem.} 2020, DOI: 10.1002/anie.202003064), metadynamics (PBE-D2, explicit CH₂Cl₂):
Good axle 6⁺: $\Delta$G‡ = 19.8 kcal/mol $\rightarrow$ K\_mod.
Bad axle 8⁺: $\Delta$G‡ \textgreater{} 100 kcal/mol $\rightarrow$ K\_trap.
Sub-Å methyl repositioning flips the barrier \textgreater{}5$\times$ at constant R\_⇌, T\_⋈.

Provisional degeneracy\_strength $\approx$ 0.71 (power-law / low-logarithmic boundary). $\Phi$\_c candidacy threshold met.

Full $\omega$B97X-D/def2-TZVPP scan pending (Tier III, P-16).

\begin{center}
\rule{0.5\textwidth}{0.4pt}
\end{center}

\subsubsection{P-15 $\cdot$ Algebra correctness suite --- 8 algebraic properties from primitive operations (V-6)  \{\#p-15-tier2\}}

\textbf{Primitive basis.} The SynthonM monad stack and the SynthOmnicon algebra rules are derived from primitive ordinals and lattice operations alone. Eight algebraic properties were predicted to hold across all \texttt{.syn} design programs and confirmed by the full 20-design suite:

\begin{table}[htbp]
\centering
\small
\rowcolors{1}{white}{white}
\begin{tabularx}{\textwidth}{>{\centering\arraybackslash}X >{\centering\arraybackslash}X >{\centering\arraybackslash}X}
\toprule
\rowcolor{gray!25}\textbf{Property} & \textbf{Mechanism} & \textbf{Designs} \\
\midrule
\rowcolors{1}{white}{gray!10}
F-floor gate & \texttt{lift(critical)} blocked when F \textless{} $F_{\hbar}$ & 01, 04 (intentional) \\
Tensor bottleneck & F = min(F₁, F₂) & 12, 16 \\
Topology promotion & T\_cage \otimes T\_cage $\rightarrow$ T\_cage; T\_⋈ \otimes T\_cage $\rightarrow$ T\_cage & 04b, 07 \\
MI discount ordering & P\_minus \otimes P\_plus $\rightarrow$ highest MI (complementary charges) & 10, 16 \\
mplus recovery & BLOCKED $\rightarrow$ join fallback $\rightarrow$ lift & 14 \\
Axiom 6 propagation & D\_$\infty$ \otimes D\_△ $\rightarrow$ D\_△$\infty$; temporal grounding carries through & 04b \\
Factor 7 operationalised & D\_$\infty$ + T\_⋈ + P\_DA + $F_{\hbar}$ $\rightarrow$ phi\_c\_score \textgreater{} 0.3 & 01b, 06, 11 \\
Path cost & Same-cluster path = 0.000 nat; cross-cluster = 0.962 nat & 06, 08, 09, 13 \\
\bottomrule
\end{tabularx}
\end{table}

\textbf{18/20 designs confirmed. 2 intentional F-floor demonstrations (designs 01, 04).}

\begin{center}
\rule{0.5\textwidth}{0.4pt}
\end{center}

\subsubsection{P-20 $\cdot$ $\lambda$ is the primitive matching fraction --- no free parameter in tensor (Occam Target 1)}

\textbf{Prediction.} The mutual-information discount factor $\lambda$ in the tensor formula \xi\_ens = \xi₁ + \xi₂ $-$ $\lambda$$\cdot$I(s₁;s₂) is not a tunable constant. It is derivable from the primitive overlap fraction:

\[\lambda(s_1, s_2) = \mathrm{frac}(s_1, s_2) = \frac{\#\{\text{matching primitive slots in } \{D,T,R,P,F,K,G\}\}}{7}\]

\textbf{Derivation.} Two boundary conditions uniquely determine $\lambda$:

\begin{enumerate}
    \item \textbf{Idempotency:} s \otimes s = s. When frac=1.0 (identical tuples), I = min(\xi₁,\xi₂) and \xi\_ens = \xi₁+\xi₂$-$min(\xi) = max(\xi) = \xi₁ (when \xi₁=\xi₂). This requires $\lambda$=1 at frac=1. 
    \item \textbf{Full synergy:} frac=0 $\rightarrow$ no MI discount $\rightarrow$ \xi\_ens = \xi₁+\xi₂. This requires $\lambda$=0 at frac=0. 
\end{enumerate}

The linear interpolation $\lambda$=frac is the unique function satisfying both.

\textbf{Verification.}

\begin{itemize}
    \item design 16 (identical cage pair, frac=1.00): \xi(s\otimess, $\lambda$=1.0) = 8.5883 = \xi(s) exactly (idempotency proven). 
    \item E{[}frac{]} over all 465 catalog pairs (first 32 synthons) = \textbf{0.3023}, within 0.002 of $\lambda$\_fixed=0.30.
    \item The fixed $\lambda$=0.30 is the catalog-mean approximation of the derived variable.
\end{itemize}

\textbf{Implication.} The tensor formula has \textbf{zero free parameters} once frac is substituted for $\lambda$. The previously-fixed $\lambda$=0.30 was the expected value of the per-pair matching fraction under the current catalog distribution --- a good approximation but not the fundamental form.

\begin{center}
\rule{0.5\textwidth}{0.4pt}
\end{center}

\subsubsection{P-21 $\cdot$ F-tier boundaries are integer Boltzmann discrimination ratios (Occam Target 2)}

\textbf{Prediction.} The numerical values F\_ell=0.40, F\_eth=0.75, $F_{\hbar}$=0.95 are not empirically chosen thresholds. They are the \textbf{Boltzmann discrimination fractions} for integer selectivity ratios:

\[F_n = \sigma\!\left(\frac{\Delta\Delta G_n}{kT}\right) = \frac{1}{1 + e^{-\Delta\Delta G_n / kT}}\]

\begin{table}[htbp]
\centering
\small
\rowcolors{1}{white}{white}
\begin{tabularx}{\textwidth}{>{\centering\arraybackslash}X >{\centering\arraybackslash}X >{\centering\arraybackslash}X >{\centering\arraybackslash}X >{\centering\arraybackslash}X >{\centering\arraybackslash}X}
\toprule
\rowcolor{gray!25}\textbf{Tier} & \textbf{F value} & \textbf{Fraction} & \textbf{Selectivity ratio} & \textbf{$\Delta$$\Delta$G/kT} & \textbf{Physical regime} \\
\midrule
\rowcolors{1}{white}{gray!10}
F\_ell & 0.40 & 2/5 & 2:3 (sub-threshold) & $-$ln(3/2) = $-$0.405 & Competition wins 3:2; 1.5 natural competitors \\
F\_eth & 0.75 & 3/4 & 3:1 (classical) & +ln(3) = +1.099 & 1 cooperative classical bond, 3:1 selectivity \\
$F_{\hbar}$ & 0.95 & 19/20 & 19:1 (quantum) & +ln(19) = +2.944 & 2 coop. bonds with enhancement; $\approx$ e³/kT $\approx$ 20:1 \\
\bottomrule
\end{tabularx}
\end{table}

\textbf{Exactness.}

\begin{itemize}
    \item logit(3/4) = ln(3) exactly 
    \item logit(19/20) = ln(19) exactly 
    \item F\_ell = 1/(1+N\_comp) = 1/2.5 = 0.40 exactly, with N\_comp = 1.5 natural competing partners 
\end{itemize}

\textbf{Cooperativity derivation.} F\_eth = 1 bond with $\Delta$$\Delta$G₀ = 1.099 kT (ln 3 kT). $F_{\hbar}$ = 2 cooperative bonds: 2$\times$1.099 = 2.197 kT (independent) + 0.747 kT cooperative enhancement = 2.944 kT = ln(19) kT. The quantum enhancement (+0.747 kT per 2-bond step) is the load-bearing part that separates $F_{\hbar}$ from the independent-bond prediction (which would give 0.900 instead of 0.950).

\textbf{Implication.} F-tier boundaries have zero free parameters. The three tiers are the only ones consistent with integer selectivity ratios while spanning sub-threshold (F\textless{}0.5), classical thermal (F\textasciitilde{}0.75), and quantum-enhanced (F\textasciitilde{}0.95) regimes.

\begin{center}
\rule{0.5\textwidth}{0.4pt}
\end{center}

\subsubsection{P-22 $\cdot$ $\Omega$ is fully derivable from \{T, K, D, $\Gamma$, G\} --- not an independent primitive (Occam Target 3)}

\textbf{Prediction.} The topological index $\Omega$ is not an independent primitive. It is an algebraic function of five existing primitives: topology (T), kinetics (K), dimensionality (D), grammar ($\Gamma$), and granularity (G).

\textbf{5-rule derivation:}

\begin{lstlisting}
\Omega(s) =
  Z2_CLASS     if T = NETWORK and D = SUPRAMOLECULAR and K in {SLOW, TRAP}
  NON_ABELIAN  if T = BRAID and \Gamma = QUANTUM_AND
  Z_CLASS      if T = LINEAR and K = TRAP and \Gamma = QUANTUM_AND
  TRIVIAL      if \Gamma = QUANTUM_AND   OR   (\Gamma = SPECIFIC_AND and G = GLOBAL)
  None         otherwise
\end{lstlisting}

Physical reading:

\begin{itemize}
    \item \textbf{Z2\_CLASS} (ℤ₂): Network topology + supramolecular bulk + gap/slow kinetics $\rightarrow$ topological insulator bulk-boundary correspondence. Determined by bulk invariant; grammar-independent.
    \item \textbf{NON\_ABELIAN}: Braided topology + quantum coherent grammar $\rightarrow$ non-Abelian anyons (FQH Moore-Read type).
    \item \textbf{Z\_CLASS} (ℤ): Linear chain + gap-protected (TRAP) + quantum coherence $\rightarrow$ Kitaev chain, Majorana zero modes.
    \item \textbf{TRIVIAL}: Quantum grammar (photon, spin, qubit) or elementary particle (SPECIFIC\_AND + G=GLOBAL for proton/electron) without topological protection signature.
    \item \textbf{None}: Classical systems --- no topological index applicable.
\end{itemize}

\textbf{Verification:} 0 mismatches across all 32 catalog synthons with the 5-rule derivation. Every synthon in the quantum cluster is correctly classified; every classical synthon correctly receives $\Omega$=None.

\textbf{Implication.} $\Omega$ is redundant as an independent primitive. The effective primitive tuple is reducible from 11 to \textbf{10 independent primitives}, with $\Omega$ a derived property. Equivalently: $\Omega$ encodes no information beyond what \{T, K, D, $\Gamma$, G\} already encode --- it is a convenience label for a 5-primitive conjunction.

\begin{center}
\rule{0.5\textwidth}{0.4pt}
\end{center}

\subsubsection{P-24 $\cdot$ Gravity theory SM/QG compatibility spectrum (§XVII)}

\textbf{Prediction.} Eight gravity theories encoded as synthon tuples and subjected to \texttt{meet}, \texttt{tensor}, \texttt{criticality\_lift}, and \texttt{path} operations against SM and canonical QG encodings. The G primitive partitions the full set: all $G_{\beth}$ theories have zero or one categorical SM conflicts and four+ categorical QG conflicts; all $G_{\aleph}$ theories invert this exactly. No theory achieves both SM and QG compatibility simultaneously.

\textbf{Key sub-results:}

\begin{itemize}
    \item \texttt{tensor(GR, SM)} $\to F_{\eth}$ bottleneck --- the quantum gravity problem as a single primitive.
    \item \texttt{criticality\_lift(AS)} BLOCKED at $G_{\beth}$ --- the Reuter fixed point is not $\Phi_c$ in the G/D sense.
    \item Hořava-Lifshitz: 4 categorical SM conflicts (P mismatch), 5 QG conflicts, $G_{\beth}$ blocking --- worst position in the space.
    \item Causal Set Theory: P conflict with both SM and QG --- categorically isolated by time-asymmetry.
    \item LQG: d(LQG, QG) minimal; one residual T conflict (network vs braid encodes the particle-statistics open problem).
    \item Entropic gravity canonically requires $T_{\cup}$ (bowl) --- first physical theory outside chemistry to require open-cavity topology. \texttt{analogies(Verlinde\_screen)} returns calixarene-class hosts.
    \item AdS/CFT: $d(\text{AdS/CFT}, \text{QG}) \approx 0.07$ --- closest existing approximation; residual gap encodes the background-independence problem as a hybrid D assignment.
\end{itemize}

\textbf{Status:} Tier II --- computationally/algebraically validated from primitive encodings.

\textbf{Method:} Primitive encoding + \texttt{meet}/\texttt{tensor}/\texttt{criticality\_lift}/\texttt{path} algebra. No domain equations inserted.

\textbf{Falsifiable sub-claim:} If Asymptotic Safety exhibits G/D degeneracy (e.g., via a holographic dual), \texttt{criticality\_lift(AS)} should return unblocked. If it does not, the blocking result stands and Asymptotic Safety's UV fixed point is not a $\Phi_c$ event in the framework's sense.

\begin{center}
\rule{0.5\textwidth}{0.4pt}
\end{center}

\subsection{Tier III --- Falsifiable, Pending Experimental Test}

\subsubsection{P-15b $\cdot$ K\_trap $\rightarrow$ K\_MBL energy cost (universal prediction)}

\textbf{Prediction.} Experimental preparation of an MBL phase from any gap-protected topological phase (Kitaev chain, FQH state, or spin-singlet correlated state) should require a free-energy cost consistent with:

\[\Delta\xi \approx 2.303 \text{ nats} \approx \ln(10) \text{ nats}\]

\textbf{per degree of freedom}, independent of the specific topological invariant (ℤ, ℤ\_class, non-Abelian). Measurable via thermodynamic integration, quench spectroscopy, or specific-heat anomaly at the MBL transition.

Falsifiable: if $\Delta$\xi measured at the MBL boundary differs systematically across topological classes (e.g., ℤ\_class vs. non-Abelian systems give different $\Delta$\xi), the K ordinal encoding is insufficient and a finer K tier is required.

\textbf{Framework confidence:} HIGH --- arithmetic consequence of K ordinal structure across confirmed topological catalog.

\begin{center}
\rule{0.5\textwidth}{0.4pt}
\end{center}

\subsubsection{P-23 $\cdot$ Standard Model $\leftrightarrow$ Quantum Gravity: structural disparity as a primitive mismatch}

\textbf{Encoding.} The Standard Model and a background-independent quantum gravity regime were encoded as synthon tuples using only the existing eleven primitives. No new primitives were added. All encodings are falsifiable by alternative primitive choices.

\textbf{Standard Model:}
\texttt{\langleD\_triangle; T\_network; R\_subset; P\_pm\_sym; $F_{\hbar}$; K\_fast; G\_beth; \Gamma\_sel-and; Phi\_sub; S=---; \Omega=---\rangle}
(D\_triangle = fixed supramolecular/multi-scale; T\_network = U(1)$\times$SU(2)$\times$SU(3) gauge coupling graph; R\_subset = directed gauge coupling; K\_fast = perturbative; \textbf{G\_beth = local gauge invariance} is the load-bearing encoding; $\Phi$\_sub = perturbative QFT)

\textbf{Quantum Gravity:}
\texttt{\langleD\_inf; T\_braid; R\_superset; P\_pm\_sym; $F_{\hbar}$; K\_trap; G\_ℵ; \Gamma\_q-and; Phi\_c; S=---; \Omega\_NA\rangle}
(D\_$\infty$ = emergent spacetime; T\_braid = braided spin networks; R\_superset = holographic/entanglement coupling; K\_trap = holographic code gap-protected; \textbf{G\_ℵ = bulk-boundary holographic} is the load-bearing encoding; $\Phi$\_c = spacetime emergence threshold)

\begin{center}
\rule{0.5\textwidth}{0.4pt}
\end{center}

\textbf{Sub-prediction P-23a: SM lift to $\Phi$\_c is blocked at G=LOCAL.}

\texttt{criticality\_lift(standard\_model)} returns \texttt{applicable=False} with:

\begin{quote}
\end{quote}

The specific blocking primitive is G=G\_beth (local gauge invariance). The Standard Model cannot be lifted to the criticality threshold from within its own primitive regime. To become critical, the SM would need G=G\_aleph --- a holographic/global description. This is precisely the AdS/CFT prescription: boundary theory becomes critical when it admits a bulk holographic dual (G\_beth $\rightarrow$ G\_aleph). The framework identifies \textbf{local gauge invariance itself} as the obstacle to criticality --- not any dynamical property.

Conversely, \texttt{criticality\_lift(quantum\_gravity)} returns \texttt{applicable=False} with: \emph{''Already at $\Phi$\_c --- no lift needed.''}

\textbf{Sub-prediction P-23b: Directed asymmetry --- emergence of classicality is the natural direction.}

\begin{table}[htbp]
\centering
\small
\rowcolors{1}{white}{white}
\begin{tabularx}{\textwidth}{>{\centering\arraybackslash}X >{\centering\arraybackslash}X >{\centering\arraybackslash}X}
\toprule
\rowcolor{gray!25}\textbf{Direction} & \textbf{Distance} & \textbf{Interpretation} \\
\midrule
\rowcolors{1}{white}{gray!10}
SM $\rightarrow$ QG (directed) & 8.40 nats & Crosses K gradient: K\_fast$\rightarrow$K\_trap is a DOWNGRADE in HotSwap metric \\
QG $\rightarrow$ SM (directed) & 6.90 nats & K\_trap$\rightarrow$K\_fast is an UPGRADE (free in directed metric) \\
Asymmetry & 1.217$\times$ & $\Delta$ = 1.50 nats = K weight $\times$ 3 tiers (exactly the K\_fast/K\_trap penalty) \\
\bottomrule
\end{tabularx}
\end{table}

QG $\rightarrow$ SM (emergence of a perturbative effective field theory from a gap-protected critical theory) is the thermodynamically natural direction in the relational lattice. SM $\rightarrow$ QG crosses the K gradient --- a gap-protected (K\_trap) target cannot be reached from a perturbative (K\_fast) source by incremental HotSwap.

\textbf{Sub-prediction P-23c: No path exists --- discontinuous transition required.}

\texttt{find\_path(standard\_model, quantum\_gravity, catalog)} returns \texttt{found=False}:

\begin{quote}
\end{quote}

There is no incremental path through any existing catalog synthon from SM to QG. The disparity is \textbf{categorically discontinuous} in D and T --- not merely expensive. This is the formal counterpart of the statement that ''there is no perturbative expansion that connects flat-space QFT to background-independent quantum gravity.''

\textbf{Sub-prediction P-23d: Four primitive CONFLICTS --- the four sources of the unification problem.}

Both meet(SM,QG) and join(SM,QG) flag identical conflicts:

\begin{table}[htbp]
\centering
\small
\rowcolors{1}{white}{white}
\begin{tabularx}{\textwidth}{>{\centering\arraybackslash}X >{\centering\arraybackslash}X >{\centering\arraybackslash}X >{\centering\arraybackslash}X >{\centering\arraybackslash}X}
\toprule
\rowcolor{gray!25}\textbf{Primitive} & \textbf{SM value} & \textbf{QG value} & \textbf{Conflict} & \textbf{Physical meaning} \\
\midrule
\rowcolors{1}{white}{gray!10}
D & SUPRAMOLECULAR & TEMPORAL &  & Fixed background vs emergent spacetime \\
T & NETWORK & BRAID &  & Local gauge coupling vs braided spin networks \\
R & COVALENT & NON\_COVALENT &  & Directed gauge coupling vs holographic entanglement \\
$\Gamma$ & SELECTIVE\_AND & QUANTUM\_AND &  & Gauge symmetry vs quantum entanglement grammar \\
P & SELF\_COMPLEMENTARY\_SYM & same &  & CPT $\leftrightarrow$ background independence (shared) \\
F & HIGH & same &  & Both quantum coherent (shared) \\
\bottomrule
\end{tabularx}
\end{table}

Four CONFLICT primitives, two shared. Any unifying theory must resolve all four conflicts simultaneously. The framework does not resolve them --- it identifies them precisely.

\textbf{Sub-prediction P-23e: tensor(SM, QG) forces $\Phi$=Phi\_c --- unification product is critical.}

\texttt{tensor(standard\_model, quantum\_gravity)} yields:

\begin{itemize}
    \item \textbf{$\Phi$ = Phi\_c} --- \emph{''$\Phi$: join propagates Phi\_c (criticality is join-dominant)''}
    \item \textbf{K = K\_trap} --- QG gap-traps the SM in the unification product; perturbativity is lost
    \item \textbf{G = G\_aleph} --- holographic global structure dominates; local gauge locality is absorbed
    \item \textbf{\xi\_CP = 14.02 nats} --- exceeds the entire catalog range of 6.55--8.83 nats; the unification product is \textbf{off-catalog}
    \item \textbf{frac = 2/7 = 0.286} --- only polarity and fidelity are shared; the tensor has near-maximal MI discount penalty for dissimilarity
    \item \textbf{Closest catalog synthon to tensor product: \texttt{synthon\_neutron} (d = 4.00)} --- the neutron, a composite bound state, is structurally closest to the unification product in the existing catalog
\end{itemize}

Any theory that combines SM and QG degrees of freedom must be at least critical. There is no sub-critical common ground: $\Phi$\_c dominates in \textbf{both} meet and join --- meaning no sub-critical theory can serve as the common language of SM and QG.

\textbf{Summary of P-23 results:}

\begin{table}[htbp]
\centering
\small
\rowcolors{1}{white}{white}
\begin{tabularx}{\textwidth}{>{\centering\arraybackslash}X >{\centering\arraybackslash}X >{\centering\arraybackslash}X}
\toprule
\rowcolor{gray!25}\textbf{Test} & \textbf{Result} & \textbf{Physical prediction} \\
\midrule
\rowcolors{1}{white}{gray!10}
lift(SM) & BLOCKED: G=LOCAL & Local gauge invariance prevents criticality; holographic G required \\
lift(QG) & Already at $\Phi$\_c & QG is already at the emergence threshold \\
d(SM$\rightarrow$QG) & 8.40 (directed) & Largest directed distance; crosses K gradient against natural flow \\
d(QG$\rightarrow$SM) & 6.90 (directed) & Natural direction: classicality emerges from criticality \\
find\_path & No path & D/T conflict is categorical, not continuous --- no perturbative bridge \\
meet/join & 4 CONFLICTS & D, T, R, $\Gamma$ --- the four primitive sources of the unification problem \\
tensor(SM,QG) & $\Phi$=Phi\_c, K=K\_trap, \xi=14.02 & Unification product is critical, off-catalog, and gap-trapped \\
\bottomrule
\end{tabularx}
\end{table}

\textbf{Caveats.} All results are conditional on the encoding being faithful. The SM's G=LOCAL is the most load-bearing and most defensible encoding choice. QG's K=TRAP (holographic gap) and T=BRAID (braided spin networks) are well-motivated but not uniquely determined. The framework makes these predictions sharply conditional on these choices.

\textbf{Falsifiability.} A unification theory that preserves G=LOCAL (local gauge invariance) and achieves $\Phi$\_c would falsify P-23a. A perturbative path from flat-space QFT to background-independent QG would falsify P-23c. A sub-critical common structure for SM and QG would falsify P-23e.

\textbf{Framework confidence:} MEDIUM --- the encoding choices are well-motivated but not unique. The results are structural consequences of the encoding, not of any additional physics input.

\begin{center}
\rule{0.5\textwidth}{0.4pt}
\end{center}

\subsubsection{P-25 $\cdot$ Black hole \xi\_CP = 0 at Hawking temperature; I scales with area via G\_aleph + T\_□□}

\textbf{Prediction.} For a Schwarzschild black hole encoded as $\langle D_{\triangle}; T_{\square\square}; R_{\supseteq}; P_{\pm}^{\text{sym}}; F_{\hbar}; K_{\text{trap}}; G_{\aleph}; \Gamma_{\wedge}(\text{QUANTUM}); \Phi_c; n{:}1 \rangle$:

(a) $\xi_{CP}(BH) = 0$ when $\eta_{CP}$ is evaluated at $T = T_H$ (Hawking temperature as the reference). At its own temperature, a black hole operates at perfect Landauer efficiency --- a consequence of the $S_{BH}$ cancellation in $\eta_{CP} = E_{\text{bit}}/(T_H \ln 2)$.

(b) $\eta_{CP}$ is mass-independent: $S_{BH}$ cancels from numerator and denominator. All black holes have the same constraint-propagation efficiency regardless of size --- the area law encoded as scale-invariant $\xi_{CP}$.

(c) The correct I(bits) for a $G_{\aleph}$ synthon is $I_{\text{boundary}}$ (boundary DOF), not $I_{\text{bulk}}$ (volume DOF). For a black hole with topology $T_{\square\square}$, the boundary is the horizon: $I = A/4l_P^2 / \ln 2$ bits. Bekenstein-Hawking entropy is the $G_{\aleph}$ correction to the I(bits) pipeline, geometrically specified by T.

\textbf{Required pipeline extension:} temperature-relative $\xi_{CP}$ mode ($E_{\text{bit}} = k_B T_{\text{system}} \ln 2$ at system temperature, not 298 K). G\_aleph entries require $I \to I_{\text{boundary}}$ substitution. Both are calibration extensions; the algebraic structure of the framework is unchanged.

\textbf{Falsification condition:} A black hole system whose $\xi_{CP}$, evaluated at $T_H$ with area-scaling I, is nonzero --- i.e., a physical black hole that violates the holographic bound. This is equivalent to a violation of the Bekenstein-Hawking formula.

\textbf{Status:} Tier III --- derived analytically; experimental confirmation requires a holographic system where $T_H$ is measurable and the entropy-area relation is testable (analogue gravity systems, e.g., sonic black holes with tunable Hawking temperature).

\begin{center}
\rule{0.5\textwidth}{0.4pt}
\end{center}

\subsubsection{P-26 $\cdot$ Hilbert space factorization failure recovered via P-20 idempotency boundary}

\textbf{Prediction.} \texttt{tensor(BH\_interior, BH\_exterior)} at maximal entanglement (frac $= 1$, all primitives matching between interior and exterior synthons) returns $\xi_{\text{ens}} = \max(\xi_{\text{int}}, \xi_{\text{ext}})$ rather than $\xi_{\text{int}} + \xi_{\text{ext}}$. This is the algebraic signature of Hilbert space factorization failure: combining two maximally entangled $G_{\aleph}$ systems does not add their information.

\textbf{Derivation.} From P-20: $\lambda = \text{frac} = 1$. From the tensor formula: $\xi_{\text{ens}} = \xi_1 + \xi_2 - 1 \cdot \min(\xi_1, \xi_2) = \max(\xi_1, \xi_2)$. No new axioms required --- the idempotency boundary condition of P-20 already encodes factorization failure as a continuous limit. The interpolation $\xi_{\text{ens}} = (2 - \text{frac}) \cdot \xi$ (for $\xi_1 = \xi_2 = \xi$) smoothly transitions from additive (frac $= 0$) to non-additive (frac $= 1$) as entanglement increases.

\textbf{Residual gap (Type III algebras).} The above derivation holds for Type I/II von Neumann algebras (trace-class density operator, countable states). Black hole horizons involve Type III$_1$ algebras (no trace, no density matrix, only relative entropies well-defined via Tomita-Takesaki modular theory). The framework's I(bits) pipeline must be extended to a relative-entropy form $D_{KL}(\rho \| \sigma)$ for Type III systems. This is a calibration extension, not a structural revision: the algebraic result (factorization failure at frac $= 1$) is correct; only the I computation changes.

\textbf{Falsification condition:} A pair of maximally entangled $G_{\aleph}$ synthons (frac $= 1$) for which \texttt{tensor} returns $\xi_{\text{ens}} < \max(\xi_1, \xi_2)$ --- i.e., entanglement somehow reduces the combined information below the maximum of the parts. This would violate the monotonicity of quantum mutual information and is physically excluded.

\textbf{Status:} Tier II/III --- the P-20 derivation is algebraically confirmed (Tier II); the Type III extension and experimental test in an analogue gravity system are Tier III.

\begin{center}
\rule{0.5\textwidth}{0.4pt}
\end{center}

\subsubsection{P-27 $\cdot$ ER=EPR as R-primitive degeneracy at $G_{\aleph}$; extended Axiom 5}

\textbf{Prediction.} If the ER=EPR conjecture (Maldacena-Susskind 2013) is correct, the R primitive is not fully independent at $G_{\aleph}$: $R_{\Leftrightarrow}$ (mechanical/topological --- wormhole) and $R_{\supseteq}$ (non-covalent/entanglement --- EPR pair) are physically indistinguishable in the $G_{\aleph}$ regime. The effective R dimension at $G_{\aleph}$ reduces from 4 to 3 ($R_{\Leftrightarrow} \equiv R_{\supseteq}$; $R_{\subseteq}$ and $R_{\ddagger}$ remain distinct). The effective independent primitive count at $G_{\aleph}$ is 10 $\to$ 8 (one from G/D degeneracy per Axiom 5; one from R-degeneracy per ER=EPR).

\textbf{Quantitative form.} Under the extended Axiom 5 R-degeneracy: \texttt{distance(ER\_bridge, EPR\_pair)} $= 0$ at $G_{\aleph}$ (vs current value of 0.27 with R and D treated as independent). \texttt{tensor(ER, EPR)} at full degeneracy: frac $= 1$, $\xi_{\text{ens}} = \max(\xi_{ER}, \xi_{EPR})$ --- they are the same system in dual descriptions, carrying no additional combined information.

\textbf{Falsification condition (strong form):} A physical system that simultaneously exhibits ER bridge geometry (wormhole-like topology with $R_{\Leftrightarrow}$ barrier profile --- a steric cliff rather than a smooth Morse curve) and EPR correlations ($R_{\supseteq}$ --- smooth, non-topological entanglement), where the two descriptions give measurably different $\xi_{CP}$ values at $G_{\aleph}$. This would disprove ER=EPR directly and confirm that $R_{\Leftrightarrow} \not\equiv R_{\supseteq}$ at $G_{\aleph}$.

\textbf{Falsification condition (weak form):} Any $G_{\aleph}$ system where the steric-cliff barrier profile (the operational signature of $R_{\Leftrightarrow}$) and the smooth Morse profile (the operational signature of $R_{\supseteq}$) can be independently measured and shown to be distinct. At $G_{\beth}$, this is easy (rotaxane vs H-bond). At $G_{\aleph}$, the framework predicts no instrument can distinguish them.

\textbf{The R values that do NOT degenerate at $G_{\aleph}$.} $R_{\subseteq}$ (covalent --- specific bond formation) and $R_{\ddagger}$ (catalytic --- transformation of one species by another) remain operationally distinguishable even at $G_{\aleph}$, because their distinction does not rely on the presence of a geometric background. Covalent bond formation requires specific orbital overlap regardless of background structure; catalytic recognition requires a distinction between catalyst and substrate that is preserved under holography (the boundary CFT contains both types).

\textbf{Extended Axiom 5 (proposed formal statement).} At $G_{\aleph}$ and $\Phi_c$: (i) G/D degenerate --- $G_{\aleph}$ and $D_{\infty}$ become informationally equivalent (original Axiom 5); (ii) R-degenerate --- $R_{\Leftrightarrow} \equiv R_{\supseteq}$ (ER=EPR; this paper's extension). The framework's distance function must apply a G\_aleph-conditional metric: $\delta_R(R_{\Leftrightarrow}, R_{\supseteq}) = 0$ when both synthons have $G = G_{\aleph}$; $\delta_R = 1$ otherwise (standard categorical distance). Implementation: \texttt{compute\_distance(s1, s2, g\_aleph\_mode=True)}.

\textbf{Status:} Tier III --- ER=EPR is a conjecture in quantum gravity with strong supporting evidence (Ryu-Takayanagi, firewall paradox resolution) but no proof. The framework's prediction is conditional on ER=EPR; it provides the precise primitive statement of what ER=EPR means in the synthon vocabulary.

\begin{center}
\rule{0.5\textwidth}{0.4pt}
\end{center}

\subsubsection{P-16 $\cdot$ Rotaxane $\Phi$\_c anchor (Transformation \#8, full scan)}

\textbf{Prediction.} A full constrained-relaxed scan of a DB24C8/dialkylammonium borderline-slippage pseudorotaxane at $\omega$B97X-D/def2-TZVPP + PCM(CH₂Cl₂) should find:
(i) degeneracy\_strength $\geq$ 0.70 (logarithmic class) at the steric-cliff TS
(ii) TS spatial correlation length \xi\_r $\approx$ axle-crown distance at the steric cliff (3.5--4.5 Å)
(iii) The first experimental anchor for Axiom 5 and the $\Phi$\_c HotSwap regime

If confirmed, this becomes the canonical $\Phi$\_c test case for the criticality-tolerant HotSwap predicted by Axiom 5 (a swap near the steric-cliff locus should be kinetically indistinguishable from a swap at the same F/K primitives away from the cliff).

Falsifiable: degeneracy\_strength \textless{} 0.50 would disconfirm the $\Phi$\_c candidacy and classify the system as non-critical by the framework's own metric.

\begin{center}
\rule{0.5\textwidth}{0.4pt}
\end{center}

\subsubsection{P-17 $\cdot$ Ice XIX correlation length \textless{} ice XV (neutron diffraction)}

\textbf{Prediction.} Ice XIX (\textgreater{}1.5 GPa, G\_beth assigned by the LLM without prompting) has a shorter antiferroelectric ordering correlation length than ice XV (\textasciitilde{}1.0 GPa, G\_gimel). This follows from the G primitive assignment alone: G\_beth = local, G\_gimel = mesoscale.

\textbf{Measurable by:} neutron diffraction correlation function analysis on ice XIX and XV under pressure. If correlation lengths are equal or XV \textless{} XIX, the G assignment convention requires revision.

Framework confidence: MEDIUM --- the G\_beth vs. G\_gimel assignment was made by LLM generation from the pressure-dependent O--O distance argument; the prediction is as reliable as that argument.

\begin{center}
\rule{0.5\textwidth}{0.4pt}
\end{center}

\subsubsection{P-18 $\cdot$ Universality track in tuple space (phase diagram, +2.303 nat displacement)}

\textbf{Prediction.} In a metric MDS projection of the quantum synthon catalog, adding a disordered Kitaev chain (K\_MBL, T\_\textbar{}, $\Omega$\_Z, all other primitives identical to kitaev\_chain\_majorana) should appear at a \textbf{fixed displacement of d = 2.303 nats} from kitaev\_chain\_majorana in the direction shared by all three K\_trap$\rightarrow$K\_MBL shifts. This displacement should be parallel for spin\_singlet $\rightarrow$ disordered\_spin and fqh $\rightarrow$ disordered\_fqh.

\textbf{Measurable by:} running \texttt{syncon phase-diagram} with the three MBL synthons added to the catalog and verifying the displacement direction is collinear.

This is the ''universality track'' in primitive space --- a purely syntactic prediction about catalog geometry.

\begin{center}
\rule{0.5\textwidth}{0.4pt}
\end{center}

\subsubsection{P-19 $\cdot$ \chi(T$\rightarrow$0) \textasciitilde{} T\^{}\{$-$$\gamma$\} divergence from Factor 8 (quantum criticality)}

\textbf{Prediction.} Synthons satisfying the Factor 8 fingerprint (G\_ℵ + $F_{\hbar}$ + K\_trap + $\neg$D\_$\infty$) sit at or near a quantum critical point. Near a QCP, the susceptibility \chi(T$\rightarrow$0) \textasciitilde{} T\^{}\{$-$$\gamma$\} with $\gamma$ \textgreater{} 0. The framework predicts this divergence for kitaev\_chain\_majorana and fqh\_moore\_read but not for topological\_insulator\_bi2se3 (K\_slow, not K\_trap; Factor 8 does not fire).

\textbf{Measurable by:} specific-heat or susceptibility measurements on Kitaev chain candidate materials (e.g., $\alpha$-RuCl₃, Na₂IrO₃) vs. Bi₂Se₃ topological insulator. If TI shows no divergence but Kitaev candidates do, Factor 8 is confirmed as a quantum criticality diagnostic. If TI also diverges, Factor 8 under-selects and the G\_ℵ + K\_trap combination needs refinement.

\begin{center}
\rule{0.5\textwidth}{0.4pt}
\end{center}

\subsubsection{P-28 $\cdot$ $\beta$-sheet scaffolds and active-site mutational tolerance}

\textbf{Primitive basis.} The tuple-space distance between \texttt{active\_site} ($\langle D_\wedge; T_\bowtie; R_\ddagger; P_{+-}; F_\eth; K_\text{mod}; G_\gimel; \Gamma_\wedge(\text{SPECIFIC})\rangle$) and \texttt{beta\_hairpin} ($\langle D_\wedge; T_\bowtie; R_\supseteq; P_\pm^\text{sym}; F_\hbar; K_\text{mod}; G_\beth; \Gamma_\wedge(\text{SPECIFIC})\rangle$) is $d = 2.80$ --- the smallest distance between any passive scaffold and the active-site tuple. The \texttt{alpha\_helix} ($T_{\vert}$, $\Gamma_\to$) is at $d \geq 4.50$ from active\_site. $\beta$-sheet topology ($T_\bowtie$) and specificity grammar ($\Gamma_\wedge(\text{SPECIFIC})$) are shared between scaffold and catalytic motif; helix topology ($T_{\vert}$) and sequential grammar ($\Gamma_\to$) are not.

\textbf{Prediction.} Alanine-scan mutations at active-site residues in $\beta$-sheet enzymes (TIM barrel, Rossmann fold, $\beta$-propeller) will be tolerated more frequently than equivalent mutations in helix-bundle enzymes (4-helix bundle, up-down bundle, coiled-coil) --- because the scaffold is closer to the active-site tuple and provides structural context that absorbs partial primitive disruptions. The fraction of alanine-scan mutations retaining $> 10\%$ wild-type activity should be higher in $\beta$-sheet enzyme classes.

\textbf{Falsification condition.} A systematic alanine-scan dataset (e.g., ProTherm, ProtaBank) showing no significant difference in tolerance rates between $\beta$-sheet and helix-bundle enzyme classes, matched for active-site geometry and substrate type.

\textbf{Status:} Tier II --- the distance matrix result is algebraically confirmed (protein\_tests.py); the experimental correlation requires systematic mutational data analysis.

\begin{center}
\rule{0.5\textwidth}{0.4pt}
\end{center}

\subsubsection{P-29 $\cdot$ Complex assembly bottleneck at $K_\text{slow}$ from tensor}

\textbf{Primitive basis.} \texttt{tensor(active\_site, allosteric\_domain)} reports $K_\text{slow}$ as the ensemble bottleneck primitive. The allosteric domain carries $K_\text{mod}$; the protein complex carries $K_\text{slow}$; the tensor result gives $K = \min(K_1, K_2) = K_\text{slow}$. The tensor bottleneck rule (validated as P-15, V-6 design suite: design 12, crown\otimesCB{[}n{]} F-bottleneck) predicts assembly rate is controlled by the slowest kinetic component.

\textbf{Prediction.} The rate-limiting step in heteromeric protein complex assembly is the association/conformational rearrangement of the allosteric regulatory subunit, not the catalytic subunit. Stopped-flow measurements should show $k_\text{fold} > k_\text{assemble}$ systematically for assemblies with allosteric regulatory subunits.

\textbf{Falsification condition.} Stopped-flow data showing $k_\text{assemble} \approx k_\text{fold}$ (assembly not slower than folding) for allosteric enzyme complexes --- implying the tensor bottleneck rule does not transfer from supramolecular chemistry to protein assembly.

\textbf{Status:} Tier II --- tensor bottleneck algebraically confirmed (V-6); protein-domain transfer requires stopped-flow kinetics.

\begin{center}
\rule{0.5\textwidth}{0.4pt}
\end{center}

\subsubsection{P-30 $\cdot$ Allosteric ON state: multi-timescale $R_{ex}$ spectrum from $\Phi_c$}

\textbf{Primitive basis.} The allosteric domain is the only protein structural unit satisfying the G/D degeneracy condition for $\Phi_c$ (Varma probe score 0.60). The Primitive Jacobian shows 6 of 8 primitives trigger Axiom 4 violations --- the allosteric domain is axiom-fragile by design. Near-critical systems (§IX) exhibit broad fluctuation spectra because they operate simultaneously at multiple length and time scales; a single-Lorentzian $R_{ex}$ spectrum indicates a single dominant timescale, characteristic of subcritical systems.

\textbf{Prediction.} The conformational exchange spectrum of allosteric domains in the signal-active (ON) state --- measured by CPMG NMR relaxation dispersion --- should show contributions across multiple timescales simultaneously (microsecond to millisecond $R_{ex}$ distribution). The signal-inactive (OFF) state, further from the criticality locus, should show a single dominant Lorentzian or no significant $R_{ex}$. Test system: CAP protein (cAMP-dependent allosteric transcription factor), whose ON/OFF switch is structurally characterized.

\textbf{Falsification condition.} CPMG data showing a single-Lorentzian $R_{ex}$ spectrum for the ON state of a canonical allosteric protein, indistinguishable from the OFF-state spectrum.

\textbf{Status:} Tier III --- the $\Phi_c$ assignment to allosteric domains is primitive-derived; the multi-timescale spectral prediction requires NMR measurement.

\begin{center}
\rule{0.5\textwidth}{0.4pt}
\end{center}

\subsubsection{P-31 $\cdot$ Directed distance asymmetry: rescue cheaper than aggregation (K-targeting corollary)}

\textbf{Primitive basis.} Directed tuple distances (F-floor asymmetry, protein\_tests3.py) for all three amyloid systems give rescue cost $<$ aggregation cost:

\begin{table}[htbp]
\centering
\small
\rowcolors{1}{white}{white}
\begin{tabularx}{\textwidth}{>{\centering\arraybackslash}X >{\centering\arraybackslash}X >{\centering\arraybackslash}X >{\centering\arraybackslash}X}
\toprule
\rowcolor{gray!25}\textbf{Disease} & \textbf{Aggregation cost} & \textbf{Rescue cost} & \textbf{Asymmetry} \\
\midrule
\rowcolors{1}{white}{gray!10}
A$\beta$ & 7.20 & 6.90 & 0.30 \\
Tau & 4.80 & 3.90 & 0.90 \\
$\alpha$-Syn & 4.80 & 3.90 & 0.90 \\
\bottomrule
\end{tabularx}
\end{table}

The asymmetry arises because the F-floor rule blocks upward $F$ transitions ($F_\ell \to F_\hbar$ in aggregation --- requires discontinuous topology change) while kinetic downgrade ($K_\text{trap} \to K_\text{fast}$ in rescue) is not floor-constrained.

\textbf{Prediction.} Kinetic-targeting rescue agents (seeding inhibitors, disaggregases, $K$-targeting chaperones) outperform thermodynamic denaturants ($F$-targeting compounds) in fibril dissolution assays. The differential should follow the rank order A$\beta$ $<$ tau $\approx$ $\alpha$-Syn, matching the directed-distance asymmetry.

\textbf{Falsification condition.} $F$-targeting denaturants dissolving fibrils as rapidly as $K$-targeting disaggregases at matched thermodynamic driving force, or the rank order not following the predicted asymmetry values.

\textbf{Status:} Tier II --- directed-distance asymmetry algebraically confirmed; experimental rank comparison requires head-to-head kinetics.

\begin{center}
\rule{0.5\textwidth}{0.4pt}
\end{center}

\subsubsection{P-32 $\cdot$ Allosteric-interface chimera: binding retained, Hill coefficient $n_H \to 1$}

\textbf{Primitive basis.} \texttt{meet(allosteric\_domain, protein\_complex)} produces two CONFLICTS: on $P$ ($P_{+-}$ vs. $P_\pm^\text{sym}$) and on $\Gamma$ ($\Gamma_\to(\text{SELECTIVE})$ vs. $\Gamma_\wedge(\text{SPECIFIC})$). These conflicts encode a structural incompatibility between allosteric signaling grammar (sequential, directional) and interface grammar (simultaneous, symmetric). A chimeric protein at this conflict boundary retains the common primitives ($D_{\wedge\triangle}$, $T_\in$, $R_\supseteq$, $F_\eth$, $K_\text{mod}$, $G_\gimel$) --- and therefore binding affinity --- but loses the coherent grammar required for cooperativity.

\textbf{Prediction.} A chimera between an allosteric regulatory domain and a quaternary interface domain retains binding to both cognate partners ($K_d$ unchanged within 10$\times$) but loses cooperativity: Hill coefficient $n_H \to 1.0$. ITC should show a normal enthalpic signature but linear binding isotherm; Hill analysis on a cooperative reporter should show $n_H$ collapse.

\textbf{Falsification condition.} A chimera at this conflict boundary retaining $n_H > 1.5$, inconsistent with the meet-conflict prediction of grammar incompatibility.

\textbf{Status:} Tier III --- requires chimera protein engineering and biophysical characterization.

\begin{center}
\rule{0.5\textwidth}{0.4pt}
\end{center}

\subsubsection{P-33 $\cdot$ Allosteric inhibitors: selectivity advantage from $G_\gimel$ vs. $G_\beth$}

\textbf{Primitive basis.} Orthosteric drugs bind at $R_\ddagger$, $G_\beth$ (local scale --- active-site pocket, shared across enzyme family members). Allosteric drugs bind at $R_\supseteq$, $G_\gimel$ (mesoscale --- full allosteric communication path, more divergent across paralogs than the catalytic pocket). By the $G$ primitive's scale-of-control definition: $G_\beth$ selectivity window = local pocket geometry; $G_\gimel$ selectivity window = domain-level contact network. The domain-level network is more divergent across kinase paralogs than the ATP-binding cleft.

\textbf{Prediction.} For any kinase family with a conserved active site, allosteric inhibitors targeting the regulatory domain ($G_\gimel$) will show higher selectivity (lower off-target activity against family paralogs) than orthosteric inhibitors targeting the ATP-binding site ($G_\beth$), in a matched selectivity panel (e.g., KINOMEscan, Ambit), even when both compounds have equivalent $K_i$ against the primary target.

\textbf{Falsification condition.} A selectivity panel showing equivalent or superior selectivity for an orthosteric compound vs. its allosteric analog in the same kinase family.

\textbf{Status:} Tier II --- the $G$ primitive assignment is the operative prediction; selectivity panel comparison is available in published kinase inhibitor datasets.

\begin{center}
\rule{0.5\textwidth}{0.4pt}
\end{center}

\subsubsection{P-34 $\cdot$ A$\beta$--$\alpha$-Syn cross-seeding rate $\approx$ homoseeding rate; both faster than tau}

\textbf{Primitive basis.} \texttt{tuple\_distance(Abeta\_fibril, alphaSyn\_fibril)} $= 0.00$. Identical primitive encoding across all 8 dimensions: both are $\langle D_{\wedge\triangle}; T_\in; R_\supseteq; P_\pm^\text{sym}; F_\hbar; K_\text{trap}; G_\aleph; \Gamma_\wedge(\text{BROAD})\rangle$. The same synthon is the same nucleation template --- there is no primitive distinction that would impose a kinetic penalty for heterologous seeding relative to homologous seeding. Tau PHF differs at $d = 2.90$ ($T_\bowtie$ vs. $T_\in$, $P_\pm^\psi$ vs. $P_\pm^\text{sym}$), predicting a cross-seeding barrier.

\textbf{Prediction.} Cross-seeding kinetics (A$\beta$ seeds $\rightarrow$ $\alpha$-Syn monomer; $\alpha$-Syn seeds $\rightarrow$ A$\beta$ monomer) should show nucleation lag times and elongation rates approximately equal to homoseeding controls (A$\beta$ seeds $\rightarrow$ A$\beta$ monomer; $\alpha$-Syn seeds $\rightarrow$ $\alpha$-Syn monomer). Both should be substantially faster than tau seeding into either A$\beta$ or $\alpha$-Syn monomer. Measurable by ThT fluorescence kinetics with exogenous seeds at sub-threshold concentration.

\textbf{Falsification condition.} A systematic cross-seeding experiment (matched seed concentration, monomer concentration, buffer) showing cross-seeding lag time $> 3\times$ homoseeding lag time for the A$\beta$ / $\alpha$-Syn pair.

\textbf{Status:} Tier II --- $d = 0.00$ is algebraically confirmed; the kinetic equality is consistent with published cross-seeding reports (Guo et al. 2013, \emph{Nat. Neurosci.}) but has not been systematically tested with matched homoseeding controls.

\begin{center}
\rule{0.5\textwidth}{0.4pt}
\end{center}

\subsubsection{P-35 $\cdot$ $K$-targeting fibril dissolution outperforms $F$-targeting denaturants}

\textbf{Primitive basis.} P-31 establishes the directed-distance asymmetry. Here the corollary is sharpened to a compound-class prediction: $F$-targeting agents (chaotropes, detergents) must navigate the F-floor to dissolve $F_\hbar$ fibrils, requiring a discontinuous topology change. $K$-targeting disaggregases (Hsp70+Hsp40+NEF) reduce $K_\text{trap}$ --- a path not floor-constrained. The tuple-space distance from fibril to monomer is 0.90 nat shorter in the $K$-targeted direction (P-31 table, tau and $\alpha$-Syn).

\textbf{Prediction.} In a head-to-head dissolution assay at matched thermodynamic driving force, $K$-targeting disaggregase systems dissolve preformed amyloid fibrils faster than $F$-targeting chemical denaturants. The rate differential should rank A$\beta$ $<$ tau $\approx$ $\alpha$-Syn, following the directed-distance asymmetry: A$\beta$ has the smallest asymmetry (0.30) and thus the smallest predicted advantage.

\textbf{Falsification condition.} GdnHCl or urea at sub-denaturing concentrations dissolving preformed fibrils as rapidly as a chaperone system at matched thermodynamic cost. Or: the dissolution rate differential not following the predicted A$\beta$ $<$ tau $\approx$ $\alpha$-Syn rank.

\textbf{Status:} Tier II --- directed-distance result algebraically confirmed; chaperone vs. denaturant comparison and rank ordering require experimental measurement.

\begin{center}
\rule{0.5\textwidth}{0.4pt}
\end{center}

\subsubsection{P-36 $\cdot$ Bivalent GNF-2: $T_\perp \to T_\in$ closes the gap to the ideal allosteric inhibitor}

\textbf{Primitive basis.} GNF-2 encodes $T_\perp$ (single branched pocket --- myristoyl binding site only). The ideal allosteric inhibitor encodes $T_\in$ (distributed contact network). The per-primitive gap analysis (protein\_tests4.py) identifies this $T$ mismatch as one of the two primitive gaps between GNF-2 ($d = 2.80$ from ideal) and the allosteric ideal. A bitopic inhibitor spanning the myristoyl pocket and a second regulatory site simultaneously shifts $T_\perp \to T_\in$, closing half the remaining distance.

\textbf{Prediction.} A bivalent GNF-2 analog engaging the ABL myristoyl pocket and a second allosteric site (e.g., the SH2-kinase linker contact) simultaneously should show: (i) lower $\xi_{CP}$ than GNF-2 monovalent (more Landauer-efficient constraint propagation); (ii) higher kinase selectivity (distributed contact grammar more paralog-discriminating); (iii) slower resistance emergence ($G_\gimel$ multi-site engagement tolerates single-site mutations). Design direction is analogous to published bitopic type I½ kinase inhibitors.

\textbf{Falsification condition.} A bivalent GNF-2 analog with $T_\in$ character showing lower selectivity and faster resistance emergence than the GNF-2 monovalent --- contradicting the $T_\perp \to T_\in$ primitive upgrade prediction.

\textbf{Status:} Tier III --- requires medicinal chemistry synthesis and SAR characterization.

\begin{center}
\rule{0.5\textwidth}{0.4pt}
\end{center}

\subsubsection{P-37 $\cdot$ $G_\beth$ drugs: binary resistance; $G_\gimel$ drugs: incremental resistance}

\textbf{Primitive basis.} Imatinib encodes $G_\beth$ (local scale): binding grammar operates at the DFG-pocket scale. Any pocket mutation disrupts binding without compensating alternative grammar --- the drug has zero contact redundancy. GNF-2 encodes $G_\gimel$ (mesoscale): binding grammar propagates across the allosteric domain. A single myristoyl-pocket mutation changes one node of the contact network but does not destroy full signal propagation. The $G$ primitive's scale-of-control definition directly predicts the resistance trajectory topology.

\textbf{Prediction.} In serial passage resistance evolution experiments: (i) imatinib-resistant clones require complete drug class switching --- pocket mutation abolishes all binding (verified: T315I, E255K $\rightarrow$ imatinib $\to$ ponatinib/asciminib); (ii) GNF-2-resistant clones show incremental resistance --- the first resistance mutation reduces, but does not abolish, sensitivity. The first passage mutation in a GNF-2 serial passage experiment should produce partial resistance (e.g., $\text{IC}_{50}$ shift 5--20$\times$), not the binary switch ($> 1000\times$ shift) characteristic of $G_\beth$ drugs.

\textbf{Falsification condition.} GNF-2 showing the same binary resistance profile as imatinib --- a first mutation abolishing all detectable inhibition --- contradicting the $G_\gimel$ multi-contact tolerance prediction.

\textbf{Status:} Tier II --- $G_\beth$ irreversibility is confirmed by clinical imatinib resistance data; the $G_\gimel$ incremental-resistance prediction for GNF-2 is a derived consequence testable by published serial passage protocols.

\begin{center}
\rule{0.5\textwidth}{0.4pt}
\end{center}

\subsection{Programmable Matter Domain (§XIX) --- P-38 through P-47}

\emph{Derived from \texttt{programmable\_matter\_tests1.py} (11 synthon encodings, pairwise distance matrix, meet operations, path algebra, Varma probe, cross-domain analogy) and \texttt{programmable\_matter\_tests2.py} (Primitive Jacobian, tensor products, DesignPipeline monad). All predictions algebraically confirmed; all await experimental validation.}

\begin{center}
\rule{0.5\textwidth}{0.4pt}
\end{center}

\subsubsection{P-38 --- Programmability pair distance predicts switching energy rank (§XIX)}

\textbf{Primitive basis.} The tuple distance between the two states of a programmable material (its ''programmability pair distance'') measures how many primitives change during a state transition and by how much. This distance should correlate with the switching energy --- the thermodynamic cost of crossing from one state to the other --- because each primitive change corresponds to a real physical difference (F = binding hierarchy, K = kinetic barrier, G = scale of order). Three material classes have cleanly resolved pair distances: $d_{\text{SMP}} = 1.70$, $d_{\text{LC}} = 3.10$, $d_{\text{colloidal}} = 5.10$.

\textbf{Prediction.} The switching energy (measurable as enthalpy: Tg-crossing enthalpy for SMP, N$\rightarrow$I transition enthalpy for LC, melting enthalpy of colloidal crystal) follows the same rank: $\Delta H_{\text{SMP}} < \Delta H_{\text{LC}} < \Delta H_{\text{colloidal crystal}}$. No domain-specific thermodynamic equation is inserted between the distance rank and the energy rank prediction --- the ordinal relationship is derived purely from the primitive metric.

\textbf{Falsification condition.} Any experimental rank inversion of switching enthalpies across these three material classes.

\textbf{Status:} Tier II --- algebraically derived from primitive distance matrix; experimental calorimetry values known in literature (DSC for SMP: \textasciitilde{}20--40 J/g; LC: \textasciitilde{}2--10 J/g; colloidal: \textasciitilde{}100+ J/g per particle contact) consistent with prediction but not used as inputs.

\begin{center}
\rule{0.5\textwidth}{0.4pt}
\end{center}

\subsubsection{P-39 --- Condensate gel rescue requires $\Gamma$-targeting or K-targeting; thermal stimulus alone fails (§XIX)}

\textbf{Primitive basis.} Condensate gel is encoded as $\langle D_{\triangle\wedge}; T_\text{network}; R_\superset; P_{\text{pm\_sym}}; F_\hbar; K_\text{trap}; G_\aleph; \Gamma_\wedge(\text{BROAD}); \Phi_{\text{sub}} \rangle$. The path search from condensate\_gel to condensate\_liquid is blocked: the F-floor theorem prohibits HotSwap traversal when $F_{\text{src}} = F_\hbar > F_{\text{dst}} = F_\ell$ at $K_\text{trap}$. No intermediate synthon in the catalog provides a stepping stone. Thermal perturbation acts only on $K$ --- it cannot lower $F$ alone.

\textbf{Prediction.} Gel dissolution requires: (a) a competing binder that lowers $F$ ($\Gamma$-targeting agent --- sequestrates the gel-forming interaction), or (b) a disaggregase/chaperone that lowers $K_\text{trap}$ (K-targeting). A stimulus that acts only on temperature (pure thermal) cannot rescue $K_\text{trap}$ gels at physiological temperatures, because temperature operates on $K$ and cannot by itself change the $F$-ordering of competing interactions.

\textbf{Falsification condition.} A purely thermal dissolution of a condensate gel (no competing ligand, no chaperone activity) at temperatures compatible with cell viability.

\textbf{Status:} Tier II --- path blocked algebraically; consistent with known biology (Hsp70 disaggregases required for TDP-43/FUS gel rescue; thermal dissolution alone requires denaturing temperatures).

\begin{center}
\rule{0.5\textwidth}{0.4pt}
\end{center}

\subsubsection{P-40 --- T-conflict programmability pairs predict first-order transitions (§XIX)}

\textbf{Primitive basis.} The meet operation on a programmability pair identifies which primitives conflict --- primitives that differ between the two states and cannot be reconciled by a shared lower bound. Topology ($T$) conflicts appear in: LC nematic ($T_\text{linear}$) vs isotropic ($T_\text{network}$), colloidal crystal ($T_{\text{network\_sym}}$) vs fluid ($T_\text{network}$). The HotSwap path is blocked for all $T$-conflicting pairs: ''HotSwap requires exact D and T match.'' Topology is discontinuous --- it cannot be traversed by continuous deformation; any topology change requires a discontinuous jump.

\textbf{Prediction.} All programmability pairs that show $T$-conflict in their meet operation undergo first-order (discontinuous) transitions: they exhibit latent heat, coexistence regions, and hysteresis. Pairs with no $T$-conflict (SMP: $T_\text{network}$ $\leftrightarrow$ $T_\text{network}$, condensate: $T_\text{network}$ $\leftrightarrow$ $T_\text{network}$) may be second-order or weakly first-order.

\textbf{Falsification condition.} A programmability pair with $T$-conflict in its meet that undergoes a continuous (second-order) transition with no latent heat and no coexistence.

\textbf{Status:} Tier II --- algebraically derived; LC N$\rightarrow$I is known to be weakly first-order (latent heat \textasciitilde{}2 kJ/mol), colloidal melting is first-order --- both consistent with $T$-conflict prediction.

\begin{center}
\rule{0.5\textwidth}{0.4pt}
\end{center}

\subsubsection{P-41 --- Actin-DNA hybrid composites show collective motion only with ATP (§XIX)}

\textbf{Primitive basis.} Tensor product: $\text{active\_gel} \otimes \text{dna\_strand\_disp}$. The composite acquires $\Phi_c$ from the active gel component (criticality is join-dominant) and $G_\aleph$ (global scale from active gel), $R_\superset$ (non-covalent recognition from DNA, downgrading the DYNAMIC\_CATALYTIC mode). The $\Phi_c$ assignment is conditional on the active gel primitive, which requires $D_{\triangle\infty}$ (temporal cycle = ATP hydrolysis). Removing ATP collapses $D_{\triangle\infty} \to D_\triangle$ (no temporal cycle), which changes the composite to a passive DNA network without $\Phi_c$.

\textbf{Prediction.} Actin-DNA composite networks (e.g., kinesin-DNA nanostructures, actin-DNA scaffolds): (1) show collective motion and spatial order absent in either component alone --- emergent from tensor composite; (2) lose collective motion upon ATP depletion faster than DNA scaffold degrades; (3) the threshold ATP concentration for collective motion onset matches Varma criticality score $\approx$ 0.60 (onset near, but below, full criticality).

\textbf{Falsification condition.} ATP-depleted actin-DNA network shows equivalent spatial order (correlation length, velocity correlations) to ATP-active network.

\textbf{Status:} Tier III --- composite derived from tensor algebra; actin-DNA composites are known but the specific ATP-dependency of collective motion onset has not been tested with this framing.

\begin{center}
\rule{0.5\textwidth}{0.4pt}
\end{center}

\subsubsection{P-42 --- Maximally versatile programmable matter is generically near-critical (§XIX)}

\textbf{Primitive basis.} The meet (dynamic floor) of all fluid/dynamic programmable matter states carries $\Phi_c$: $\langle D_\triangle; T_\text{network}; R_\superset; P_{\text{pm\_sym}}; F_\ell; K_\text{mod}; G_\beth; \Gamma_\vee(\text{BROAD}); \Phi_c \rangle$. This is not a design choice inserted into any individual encoding --- it emerges from the lattice meet of six independently encoded systems. $\Phi_c$ is join-dominant in the meet lattice: it propagates whenever at least one component is near-critical.

\textbf{Prediction.} Any programmable matter system engineered to maximize the number of accessible states (highest programmability) will converge to near-critical encoding: $F_\ell$, $K_\text{mod}$, $G_\text{global}$, $\Phi_c$. This is not because designers choose criticality --- it is because the algebra shows $\Phi_c$ is the only phase assignment compatible with the dynamic floor. Consequence: the most versatile programmable matter in biology (cytoplasm, membraneless organelles) operates near criticality --- consistent with measured cortical criticality, now derived from primitive algebra without biological inputs.

\textbf{Falsification condition.} A maximally versatile programmable matter system (demonstrated large state space, global responsiveness) that is provably subcritical in all observables.

\textbf{Status:} Tier II --- lattice algebra result; consistent with growing evidence for criticality in biological PM systems (cortex, cytoplasm); prediction extends to synthetic PM design.

\begin{center}
\rule{0.5\textwidth}{0.4pt}
\end{center}

\subsubsection{P-43 --- Condensates implement mesoscale allostery with Hill coefficient tied to criticality score (§XIX)}

\textbf{Primitive basis.} Cross-domain analogy: nearest catalog neighbor to condensate\_liquid is allosteric\_domain at $d = 2.50$ (closer than any molecular synthon). The shared primitives are: $F_\ell$ (individually weak contacts), $K_\text{fast}$ (dynamic exchange), $G_\gimel$ (mesoscale scale of control), $\Gamma_\vee$ (promiscuous partner), $\Phi_c$ (near-critical). This is the same tuple structure that defines a mesoscale allosteric system: many weak contacts at intermediate scale, dynamic, near-critical, promiscuous.

\textbf{Prediction.} (1) Condensate-mediated signaling shows distance-independent propagation within the droplet --- a perturbation at one face propagates to the opposite face faster than diffusion-limited transport, because $G_\gimel$ + $\Phi_c$ enables global-from-local propagation. (2) The apparent Hill coefficient of condensate-mediated catalysis equals the Varma criticality score to within 15\%: $n_H \approx 0.60 / 0.50 \approx 1.2$. (3) Condensate dissolution abolishes cooperativity --- $n_H$ drops to 1.0.

\textbf{Falsification condition.} Condensate-mediated signaling shows purely diffusion-limited propagation with no superdiffusive component. Or: $n_H$ in a condensate-mediated reaction is $> 2.0$ or $< 1.0$.

\textbf{Status:} Tier III --- cross-domain analogy; the Hill coefficient prediction is novel and experimentally accessible by comparing condensate vs. dilute-phase catalysis rates.

\begin{center}
\rule{0.5\textwidth}{0.4pt}
\end{center}

\subsubsection{P-44 --- Colloidal crystals with correct symmetry support topologically protected boundary states (§XIX)}

\textbf{Primitive basis.} Nearest catalog neighbor to colloidal\_crystal is topological\_insulator\_bi2se3 at $d = 2.80$, and synthon\_neutronium at $d = 2.80$. The shared structure: $F_\hbar$ (high fidelity collective binding), $G_\aleph$ (global crystalline order), $T_\text{network}$ (lattice), $R_\superset$ (non-covalent). This is the same primitive cluster as a topological insulator. The distance $d = 2.80$ is close enough to predict structural analogy.

\textbf{Prediction.} Colloidal crystals with $d < 3.0$ from topological\_insulator in the primitive metric support topologically protected boundary modes --- surface states robust to bulk disorder --- when the interaction design matrix has the correct band-crossing symmetry (non-trivial Zak phase). The range of colloidal systems predicted to show topological boundary states is defined by $d < 3.0$: any colloidal crystal with $F_\hbar$, $G_\aleph$, $T_\text{network}$ and tuned interaction symmetry qualifies.

\textbf{Falsification condition.} A colloidal crystal with $d < 3.0$ from topological\_insulator but no detectable boundary-protected states, even after interaction symmetry tuning.

\textbf{Status:} Tier II --- phenomenon confirmed experimentally for DNA-coated colloidal crystals (Rechtsman 2016 and subsequent); prediction extends to: the primitive distance threshold ($d < 3.0$) as the boundary of the topological analogy.

\begin{center}
\rule{0.5\textwidth}{0.4pt}
\end{center}

\subsubsection{P-45 --- Topologically closed DNA origami (knots, catenanes) shows polynomial strand-failure sensitivity (§XIX)}

\textbf{Primitive basis.} DNA origami is nearest to topological\_insulator\_bi2se3 at $d = 3.70$ --- more distant than colloidal crystal ($d = 2.80$) but still within topological analogy range. DNA origami with $\Omega \neq 0$ (topological closure --- knots, catenanes, Borromean rings) acquires topological protection. In the framework, $\Omega$ measures topological robustness: how many strand-failure events are required to destroy structural integrity changes from $n$ (linear: one staple failure is irreversible) to $O(n^2)$ or better (topological: multiple simultaneous failures required).

\textbf{Prediction.} A DNA origami structure with topological closure ($\Omega \neq 0$) shows error-rate scaling as $\sim \exp(-n_\text{strands})$ only in open structures. In topologically closed structures, error scaling becomes $\sim n_\text{strands}^{-k}$ where $k$ depends on the topological class. Measurable: compare FRET-detected folding yield for open vs. catenated origami as a function of staple strand concentration.

\textbf{Falsification condition.} Topologically closed DNA origami (knots, catenanes) shows identical strand-failure sensitivity to equivalent open origami.

\textbf{Status:} Tier III --- topological DNA origami exists experimentally (Lim et al. 2020) but polynomial vs. exponential error scaling has not been tested as a function of topological class.

\begin{center}
\rule{0.5\textwidth}{0.4pt}
\end{center}

\subsubsection{P-46 --- Primitive Jacobian identifies dominant design lever: F controls \textgreater{} 40\% of d-reduction (§XIX)}

\textbf{Primitive basis.} The primitive Jacobian $\partial d / \partial \text{primitive}$ measures the sensitivity of programmability pair distance to single-primitive perturbation. Computed for all five PM pairs: F (fidelity) gives the largest $|\Delta d|$ in DNA ($-$1.10), colloidal ($-$1.20), SMP ($-$0.60), LC ($-$0.60) pairs. K (kinetic character) dominates condensate gel rescue ($-$1.50). G contributes secondary leverage in DNA ($-$0.90) and colloidal ($-$0.80) pairs but is zero for SMP (G\_gimel invariant in both states).

\textbf{Prediction.} Engineering the Jacobian-identified primitive alone achieves $> 40\%$ of the maximum possible pair distance reduction (full 11-primitive optimization). Corollary: multi-primitive optimization gives diminishing returns beyond the top two Jacobian-ranked primitives (law of diminishing returns in primitive space). Specific experimental test: vary Tg (= $F$ proxy for SMP) while holding crosslink density ($K$ proxy) and domain size ($G$ proxy) constant --- this alone should halve the switching energy relative to varying crosslink density alone.

\textbf{Falsification condition.} For any PM pair, optimizing the bottom-ranked Jacobian primitive gives equal or larger d-reduction than the top-ranked primitive.

\textbf{Status:} Tier III --- Jacobian computed algebraically; experimental Tg vs crosslink-density vs domain-size optimization comparison is feasible but not yet published with this framing.

\begin{center}
\rule{0.5\textwidth}{0.4pt}
\end{center}

\subsubsection{P-47 --- Actin-DNA composites cannot achieve a shared locked state: structural integrity requires ATP (§XIX)}

\textbf{Primitive basis.} meet(active\_gel, dna\_origami) conflicts on $\{D, R, P, \Gamma\}$. Active gel has $D_{\triangle\infty}$ (temporal cycle = ATP hydrolysis); DNA origami has $D_{\triangle\wedge}$ (molecular + supramolecular, no temporal cycle). These are incompatible at the level of dimensionality: no shared $D$ value exists that satisfies both. Materials whose rigid states conflict on $D$ have no shared locked state in the algebra --- there is no primitive tuple that is both ''static DNA scaffold'' and ''static actin network'' simultaneously.

\textbf{Prediction.} Actin-DNA composite materials cannot maintain structural integrity without ATP. The ''static DNA scaffold plus actin'' design is algebraically incoherent --- actin's structural integrity is encoded in $D_{\triangle\infty}$ (the temporal ATP cycle), so removing ATP collapses the actin component. Specific prediction: actin-DNA hybrids degrade without ATP replenishment faster than equivalent ATP-free DNA origami scaffolds, because actin depolymerisation is part of the actin synthon's $D$ primitive (not an external degradation pathway).

\textbf{Falsification condition.} An actin-DNA composite that maintains structural integrity for $> 24$ h without ATP, matching performance of equivalent ATP-free DNA origami scaffold under identical conditions.

\textbf{Status:} Tier III --- meet conflict derived algebraically; consistent with known actin biology (G-actin/F-actin equilibrium requires ATP turnover) but specific comparison with DNA origami stability has not been published.

\begin{center}
\rule{0.5\textwidth}{0.4pt}
\end{center}

\subsection{P-48 --- Condensate Gelation and Amyloid Fibrillization Are the Same Primitive Event; K-Targeting Is the Preferred Therapeutic Strategy for Both}

\textbf{Source.} Programmable matter encoding (PROGRAMMABLE\_MATTER.md §5); protein domain catalog (PROTEIN\_APPLICATIONS.md).

\textbf{Algebraic basis.} $d(\text{condensate\_gel}, \text{amyloid}) = 0.00$ --- the programmable matter gel state and the amyloid fibril are identical in all nine primitives:

\[\langle D_\triangle; T_\in; R_{\superset}; P_{\pm\text{pseudo}}; F_\hbar; K_\text{trap}; G_\gimel; \Gamma_\text{and}(\text{SELECTIVE}); \Phi_\text{sub} \rangle\]

This is not a structural analogy. A distance of zero means the primitive encoding is the same object. The two literatures (condensate biology and amyloid neuroscience) have independently identified the same physical attractor: a trapped, high-fidelity, mesoscale, supramolecular network with self-complementary polarity. The difference is biological context, not physical structure.

\textbf{Primitive Jacobian result.} The Jacobian computed over the full programmable matter catalog assigns the K primitive the highest single-primitive leverage for gel state rescue: targeting $K$ ($K_\text{trap} \to K_\text{fast}$) reduces the tuple distance from the locked state to the dynamic floor by more than targeting $F$ alone. This extends directly to amyloid: $K$-targeting (kinetic remodelling of the trapped state --- disaggregases, chaperones, ATP-dependent unfolding machines) should outperform $F$-targeting (binding competitors that lower association affinity) as a dissolution strategy.

\textbf{Prediction 1 --- Mechanistic equivalence.} A dissolution agent that successfully converts a condensate gel to the liquid state (reducing $K_\text{trap} \to K_\text{fast}$) will show non-trivial cross-reactivity against amyloid fibrils of the same polarity class ($P_{\pm\text{pseudo}}$: A$\beta$, $\alpha$-Syn, tau), because the primitive target is identical. Agents that act through the $F$ primitive (competitive binding, affinity-based disruption) will not show this cross-reactivity --- they address different positions in the tuple.

\textbf{Prediction 2 --- Jacobian hierarchy.} For any disease system where both condensate gelation and amyloid formation are co-pathological (e.g., TDP-43 in ALS, FUS in FTLD): therapeutic strategies that target K (Hsp70/Hsp104 chaperone axis, disaggregase activity, ATP-dependent remodelling) will outperform strategies that target F (small-molecule binding competitors, antibodies that raise the association barrier) in dissolution efficiency. Ratio of K-strategy to F-strategy dissolution rates should exceed 1.5$\times$ from the Jacobian prediction.

\textbf{Prediction 3 --- Locking asymmetry preservation.} The same directed-distance asymmetry that makes programmable matter locking thermodynamically downhill applies to amyloid: the forward rate (monomer $\rightarrow$ fibril, $K_\text{mod} \to K_\text{trap}$) will always exceed the reverse rate (fibril $\rightarrow$ monomer, $K_\text{trap} \to K_\text{mod}$) under identical conditions, regardless of specific amino acid sequence. The asymmetry is structural, not sequence-specific. Falsified by: any amyloid system where spontaneous dissolution exceeds nucleated aggregation at the same monomer concentration and temperature.

\textbf{Connection to P-34 and P-35.} P-34 established that A$\beta$ and $\alpha$-Syn share $d = 0.00$ (cross-seeding $\approx$ homoseeding). P-35 established that $K$-targeting dissolution ranks below $F$-targeting for the amyloid family. P-48 now links the condensate gelation literature to both: the three predictions form a coherent triangle --- zero-distance identity, K-strategy superiority, locking asymmetry --- all derivable from the same primitive tuple.

\textbf{Falsification condition.} (1) A dissolution agent that targets $F$ (not $K$) showing equal or greater cross-reactivity against both condensate gels and amyloid fibrils would falsify Prediction 1. (2) A K-strategy dissolution agent showing no advantage over F-strategy in a head-to-head assay for both target classes would falsify Prediction 2.

\textbf{Status:} Tier II --- zero-distance identity algebraically confirmed; Jacobian hierarchy computed from catalog; consistent with known biology (ATP-dependent disaggregases dissolve both condensates and amyloid; LLPS-to-gel transitions share kinetics with early amyloid nucleation). Experimental head-to-head comparison of K-targeting vs F-targeting agents against matched condensate gel and amyloid fibril targets not yet published.

\begin{center}
\rule{0.5\textwidth}{0.4pt}
\end{center}

\subsection{P-49 --- Phi\_c Categorical Independence: Critical Phase Is a Label, Not a Derived Ordinal}

\textbf{Origin.} Decomposition algebra exploration, 2026-03-19. Principal result: \texttt{kernel(allosteric\_domain, phi\_c\_probe)} strips $F$, $K$, $G$ to their floors, but the Phi=CRITICAL field survives because it is a \emph{categorical} field in the synthon dataclass --- not computed from ordinals. Constructing \texttt{phi\_c\_skeleton} ($F$=LOW, $K$=FAST, $G$=LOCAL, Phi=CRITICAL) is a valid synthon. Catalog proof: \texttt{asymptotic\_safety\_reuter\_fp} carries $G$=LOCAL + Phi=CRITICAL --- the UV fixed point of asymptotic safety is \emph{locally} critical without global organisation.

\textbf{Prediction.} Criticality-organising effects should be observable in systems with individually weak interactions ($F$=LOW, each contact $\sim 1$--$3\,k_BT$) if the network topology ($T$=NETWORK) is correct. Predicted systems: LLPS condensates (weak IDR contacts but near-critical droplet), neuronal avalanches (low synaptic $F$, global $\Phi_c$), and low-affinity multivalent receptors at crowded membranes. The Phi\_c field encodes the \emph{class} of dynamical regime --- not the magnitude of any ordinal.

\textbf{Falsification condition.} A system with $F$=LOW, $K$=FAST, $G$=LOCAL that cannot in principle support a phase transition would falsify categorical independence. (Note: the prediction is about the encoding grammar --- it does not assert that every $F$=LOW system \emph{is} critical, only that criticality is not \emph{forbidden} by low ordinals.)

\textbf{Status:} Tier II --- algebraically exact; catalog proof exists (\texttt{asymptotic\_safety\_reuter\_fp}); consistent with LLPS phenomenology. Systematic ordinal-vs-criticality screen not yet performed.

\begin{center}
\rule{0.5\textwidth}{0.4pt}
\end{center}

\subsection{P-50 --- Amyloid Formation Algebraically Requires an External High-Fidelity Nucleation Seed}

\textbf{Origin.} \texttt{cofactor(amyloid\_fibril, condensate\_liquid)} $\rightarrow$ F-CONFLICT: tensor-min($F$=MEDIUM, $B$) $\leq$ MEDIUM $<$ HIGH. A liquid condensate ($F$=MEDIUM) cannot template amyloid ($F$=HIGH) by primitive tensor alone. This is not a kinetic barrier --- it is a structural impossibility in the primitive algebra: no value of $B$ satisfies min(MEDIUM, $B$) = HIGH.

\textbf{Prediction.} Condensate-to-amyloid conversion must always be gated by an external high-fidelity nucleation event: a seed fibril, a metal ion interface (which elevates local $F$), a lipid membrane surface, or any other $F$=HIGH input. Unseeded liquid condensates of the same protein should show an indefinite lag phase that collapses upon addition of pre-formed fibril seeds. The lag phase duration should be inversely proportional to seed concentration with a saturation plateau at the $F$-floor threshold.

\textbf{Falsification condition.} A purified liquid condensate that converts to amyloid \emph{de novo} with no detectable nucleation lag, in a system verified to contain no fibril seeds or $F$=HIGH surfaces, would falsify this prediction.

\textbf{Connection to P-34, P-35, P-48.} P-34 established $d(\text{A\beta}, \text{\alpha-Syn}) = 0.00$; P-48 established condensate gel and amyloid are the same primitive event. P-50 now algebraically derives the nucleation requirement, completing the triangle: identity (P-34) + K-strategy superiority (P-48) + F-floor nucleation barrier (P-50).

\textbf{Status:} Tier II --- algebraically exact (F-CONFLICT derived); consistent with seeded aggregation literature (Jarrett \& Lansbury 1993, Bhak et al. 2009). Matched-control seeded vs. unseeded condensate kinetics under controlled seed-free conditions not yet published as a direct test.

\begin{center}
\rule{0.5\textwidth}{0.4pt}
\end{center}

\subsection{P-51 --- The Quantization Residual: cofactor(QG, GR) Is the Quantization Operator as a Synthon}

\textbf{Origin.} \texttt{cofactor(quantum\_gravity, general\_relativity)} computes residual $B$ such that GR $\otimes$ $B$ $\approx$ QG. The per-primitive analysis from \texttt{decompose\_explorations2.py} gives:

\begin{table}[htbp]
\centering
\small
\rowcolors{1}{white}{white}
\begin{tabularx}{\textwidth}{>{\centering\arraybackslash}X >{\centering\arraybackslash}X >{\centering\arraybackslash}X >{\centering\arraybackslash}X >{\centering\arraybackslash}X}
\toprule
\rowcolor{gray!25}\textbf{Primitive} & \textbf{Role} & \textbf{GR value} & \textbf{QG value} & \textbf{Interpretation} \\
\midrule
\rowcolors{1}{white}{gray!10}
$K$ & Contributor & SLOW & TRAP & Path-integral measure freezes high-action histories \\
$G$ & Contributor & LOCAL & GLOBAL & Quantisation promotes local GR to global entanglement \\
$T$ & Contributor & NETWORK & BRAID & Smooth spacetime $\rightarrow$ spin-foam / anyonic exchange topology \\
Phi & Contributor & SUBCRITICAL & CRITICAL & Quantum phase transitions have no classical analogue \\
Omega & Contributor & None & NON\_ABELIAN & Non-Abelian gauge structure / anyonic statistics \\
$D$ & \textbf{CONFLICT} & SUPRAMOLECULAR & TEMPORAL & GR organises matter in space; QG must organise spacetime --- orthogonal D-components \\
\bottomrule
\end{tabularx}
\end{table}

\textbf{Prediction.} The $D$-CONFLICT --- GR's $D$=SUPRAMOLECULAR (spatial substrate) vs. QG's $D$=TEMPORAL --- is the algebraic form of the \emph{background-dependence problem}: any quantisation scheme that preserves GR's spatial $D$-component will produce a $D$-CONFLICT at the quantum gravity level. A background-free theory of quantum gravity must either (a) abandon GR's spatial $D$-component, or (b) introduce a $D$-merging primitive not yet in the catalog. Corollary: the quantisation residual $B = \{K_\text{trap}, G_\aleph, T_\text{braid}, \Phi_c, \Omega_\text{NA}\}$ is the minimal primitive tuple that, when tensored with GR, produces the observed structure of QG --- it encodes the complete content of ''quantisation'' as a single synthon.

\textbf{Falsification condition.} A formulation of quantum gravity that (a) retains GR's spatial $D$-component and (b) does not produce a $D$-conflict with the temporal structure of quantum fluctuations would require a revision of the $D$-primitive taxonomy or the tensor semantics.

\textbf{Status:} Tier III --- algebraically derived from catalog encodings; interpretation consistent with known obstacles to canonical quantisation; not yet experimentally testable in isolation. Cross-reference: P-23 (SM/QG lift blocked), P-24 ($G$ partitions gravity theories), P-27 (ER=EPR extended Axiom 5).

\begin{center}
\rule{0.5\textwidth}{0.4pt}
\end{center}

\subsection{P-52 --- GNF-2 Drug Combination Strategy: A Supramolecular Network Scaffold Closes the Allosteric Gap}

\textbf{Origin.} Drug panel comparison via \texttt{cofactor(allosteric\_domain, drug)} for GNF-2, imatinib, and venetoclax. GNF-2 is the only drug carrying $\Phi_c$ (distance 3.5, zero conflicts). Cofactor residual $B$ = \{$D$=SUPRAMOLECULAR, $T$=NETWORK, $F$=MEDIUM, $K$=MODERATE, $G$=LOCAL, Phi=SUBCRITICAL\} --- $B$ need not be critical (GNF-2 already carries $\Phi_c$); $B$ must supply the supramolecular network topology that GNF-2 lacks.

\textbf{Prediction.} GNF-2 co-administered with a supramolecular network scaffold --- candidates include a PROTAC bifunctional linker, a DNA nanostructure template, a polyelectrolyte hydrogel matrix, or a PEGylated polyvalent nanocage --- would close the $D$+$T$ primitive gap to the allosteric domain target. The combination strategy is mechanistically distinct from simple additive inhibition: the scaffold partner does not need to bind Bcr-Abl directly; it must provide a \emph{network topology medium} ($T$=NETWORK) that allows allosteric signal propagation, while GNF-2 provides the $\Phi_c$-carrying warhead. This is the algebraic argument for a GNF-2 PROTAC or bispecific adaptor approach where GNF-2 is the warhead and the scaffold is the linker that upgrades the topology.

\textbf{Falsification condition.} If a PROTAC or scaffold co-administration of GNF-2 shows no improvement in allosteric target engagement (as measured by NMR $R_{ex}$ or FRET-based conformational reporters) compared to GNF-2 alone, the $D$+$T$ gap interpretation is incorrect.

\textbf{Connection to P-36.} P-36 predicted bivalent GNF-2 ($T_\perp \to T_\in$) lowers $\xi_{CP}$ and improves selectivity. P-52 extends this: the scaffold partner specifies the supramolecular network \emph{medium}, not just the topology class.

\textbf{Status:} Tier III --- algebraically derived; consistent with known GNF-2 allosteric mechanism (Panjarian 2013) and PROTAC feasibility. No GNF-2-scaffold co-administration study has been published as of 2026-03-19.

\begin{center}
\rule{0.5\textwidth}{0.4pt}
\end{center}

\subsection{P-53 --- Three Stability Regimes Are Algebraically Distinguishable by Peel Cost Profile}

\textbf{Origin.} Peel cost analysis of three benchmark systems from \texttt{decompose\_explorations2.py}:

\begin{table}[htbp]
\centering
\small
\rowcolors{1}{white}{white}
\begin{tabularx}{\textwidth}{>{\centering\arraybackslash}X >{\centering\arraybackslash}X >{\centering\arraybackslash}X >{\centering\arraybackslash}X >{\centering\arraybackslash}X}
\toprule
\rowcolor{gray!25}\textbf{System} & \textbf{Regime} & \textbf{Phi peel cost} & \textbf{K-floor?} & \textbf{Dissolution mechanism} \\
\midrule
\rowcolors{1}{white}{gray!10}
DNA origami & Thermodynamic & 0 & No & Gradual UV-melting; cooperative but continuous \\
Condensate liquid & Phase-protected & 3.0 nats & No & Sharp 2-state transition; Phi\_c removal discontinuous \\
Condensate gel / amyloid & Kinetically frozen & 0 & Yes ($K$=TRAP) & Mechanical disruption or specific disaggregase only \\
\bottomrule
\end{tabularx}
\end{table}

DNA origami has $\Phi$=SUBCRITICAL and topo\_index=None $\rightarrow$ zero peel costs everywhere. Condensate liquid has Phi=CRITICAL $\rightarrow$ $\Phi$ removal costs 3.0 nats; $K$ is not at floor $\rightarrow$ $K$ removal carries a finite cost. Condensate gel has $K$=TRAP already at floor $\rightarrow$ zero $K$ peel cost, but kinetically immobile; thermally inaccessible dissolution.

\textbf{Predictions.}

\begin{enumerate}
    \item LLPS condensates should show sharper melting cooperativity (steeper van't Hoff $n_H$) than DNA origami structures of equivalent thermal stability ($T_m$ matched by buffer/salt adjustment). The discontinuous 2-state cooperativity is the direct consequence of the Phi peel cost.
    \item Condensate gels and amyloid fibrils should resist thermal dissolution at temperatures that melt DNA origami of matched $T_m$. The kinetically frozen regime requires active mechanical energy input or a specific enzymatic ($K$-targeting) disaggregase.
    \item The peel cost is measurable as the enthalpy difference between the 2-state melting endpoint and the gradual baseline: $\Delta H_{\Phi_c} \approx 3.0\,k_BT \times N_{\text{contacts}}$.
\end{enumerate}

\textbf{Falsification condition.} A LLPS condensate showing the same melting cooperativity as DNA origami under identical conditions would require revising the Phi peel cost to zero, invalidating the categorical distinction.

\textbf{Status:} Tier II --- algebraically derived from peel cost accounting; consistent with known LLPS phenomenology (sharp transitions in condensate FRAP recovery; gradual DNA origami melting curves). Direct head-to-head DSC comparison of DNA origami vs. condensate with matched $T_m$ not yet published.

\begin{center}
\rule{0.5\textwidth}{0.4pt}
\end{center}

\subsection{P-54 --- AdS/CFT Holography Is the Operation GR \otimes \langleG\_GLOBAL \otimes Phi\_c\rangle}

\textbf{Origin.} \texttt{cofactor(ads\_cft\_boundary, general\_relativity)} per-primitive analysis:

\begin{table}[htbp]
\centering
\small
\rowcolors{1}{white}{white}
\begin{tabularx}{\textwidth}{>{\centering\arraybackslash}X >{\centering\arraybackslash}X >{\centering\arraybackslash}X >{\centering\arraybackslash}X}
\toprule
\rowcolor{gray!25}\textbf{Primitive} & \textbf{Role} & \textbf{GR} & \textbf{AdS/CFT} \\
\midrule
\rowcolors{1}{white}{gray!10}
$G$ & Contributor & LOCAL & GLOBAL \\
Phi & Contributor & SUBCRITICAL & CRITICAL \\
$K$ & CONFLICT & SLOW & MODERATE \\
$D$ & CONFLICT & SUPRAMOLECULAR & HOLOGRAPHIC \\
$F$ & Bottleneck & HIGH & MEDIUM \\
\bottomrule
\end{tabularx}
\end{table}

AdS/CFT contributes exactly $G$=GLOBAL and Phi=CRITICAL above GR. The cofactor $B = \langle G_\aleph \otimes \Phi_c \rangle$.

\textbf{Interpretation.}

\begin{itemize}
    \item $G$: LOCAL $\rightarrow$ GLOBAL encodes the \emph{bulk-to-boundary global correlation structure}. GR is a local theory (local diffeomorphisms). AdS/CFT makes it global: the bulk is dual to the boundary at a different scale entirely. This is the holographic renormalisation group, algebraically.
    \item Phi: SUBCRITICAL $\rightarrow$ CRITICAL encodes the \emph{boundary CFT criticality}. GR's bulk has no intrinsic criticality. AdS/CFT's boundary is a conformal field theory at a critical point. The boundary lives at $\Phi_c$ by definition.
    \item $D$-CONFLICT (SUPRAMOLECULAR vs. HOLOGRAPHIC) and $K$-CONFLICT encode the fact that the holographic dictionary is not a simple tensor product of GR --- it requires a genuinely new dimensionality primitive ($D_\text{holo}$) and a kinetic restructuring.
\end{itemize}

\textbf{Prediction.} The cofactor $B = \langle G_\aleph \otimes \Phi_c \rangle$ is the \emph{minimal primitive content of holographic duality}. Any physical mechanism that independently adds both global-scale correlation ($G$=GLOBAL) and a critical boundary ($\Phi_c$) to a local classical gravity theory should exhibit holographic-like features (bulk-boundary correspondence, entropy area law, emergent gauge symmetry). Candidate systems: analogue gravity condensates at critical points, acoustic black holes in critical quantum fluids.

\textbf{Connection to P-27 (ER=EPR).} The ER=EPR prediction required extending Axiom 5 to include $R$-degeneracy at $G_\aleph$. P-54 now shows that AdS/CFT additionally requires $G$=GLOBAL above GR --- the two predictions are consistent: both ER=EPR and AdS/CFT live in the $G_\aleph$ regime, and the distinction between them is the presence ($\Phi_c$, holography) or absence (ER=EPR geometry alone) of boundary criticality.

\textbf{Falsification condition.} An analogue gravity system that achieves $G$=GLOBAL correlation + boundary criticality without exhibiting any holographic-like entropy scaling or bulk-boundary correspondence would falsify the identification of $B = \langle G_\aleph \otimes \Phi_c \rangle$ as the holographic cofactor.

\textbf{Status:} Tier III --- algebraically derived from catalog encodings; interpretation consistent with standard AdS/CFT literature; analogue gravity test systems not yet constructed to specification.

\begin{center}
\rule{0.5\textwidth}{0.4pt}
\end{center}

\subsection{P-55 $\cdot$ Mechanical Priming Collapses to Constitutive State: d(AtHv1\_primed, PsHv1) = 0}

\textbf{Primitive basis.} AtHv1 in its electrically silent state encodes:

\[\langle D_\wedge; T_\in; R_\supset; P_{\pm\psi}; F_\eth; K_\text{trap}; G_\beth; \Gamma_\wedge(\text{SELECTIVE}); \Phi_\text{sub}; \Omega_0 \rangle\]

After mechanical priming (membrane stretch destabilizes RSN), the peel operation acts on K:

\[\text{peel}(K_\text{trap}) \Rightarrow K_\text{mod}; \quad T: T_\in \to T_\bowtie; \quad P: P_{\pm\psi} \to P_\text{directional}\]

The primed state encodes:

\[\langle D_\wedge; T_\bowtie; R_\supset; P_\text{directional}; F_\eth; K_\text{mod}; G_\beth; \Gamma_\wedge(\text{SELECTIVE}); \Phi_\text{sub}; \Omega_0 \rangle\]

PsHv1 (constitutively primed gymnosperm channel) encodes identically:

\[\langle D_\wedge; T_\bowtie; R_\supset; P_\text{directional}; F_\eth; K_\text{mod}; G_\beth; \Gamma_\wedge(\text{SELECTIVE}); \Phi_\text{sub}; \Omega_0 \rangle\]

\textbf{Algebraic result.} $d(\text{AtHv1\_primed},\ \text{PsHv1}) = 0.000$ --- exact identity in primitive space.

\textbf{Interpretation.} The RSN (Ring-Shaped Network: K117, E173, T174, K154, K155, S164) is the \emph{only primitive distinction} between AtHv1 and PsHv1. It implements K\_trap + correlated P and T changes (the channel cannot cycle or direct until K\_trap is peeled). Mechanical priming removes these three primitive differences simultaneously, yielding the gymnosperm constitutive state.

This is confirmed experimentally by the chimera transplant experiments: ChE3-4.S4.K (PsHv1 with AtHv1 KET residues) acquires mechanical sensitivity; conversely, K154-K155 + S164 substitutions in PsHv1 enhance priming. The tuple distance collapses to zero in both directions.

\textbf{Consequence.} The K\_trap peel cost (AtHv1\_silent $\rightarrow$ AtHv1\_primed) is $d = 3.3$, spanning the K, T, and P primitives simultaneously. This is the largest single-operation primitive cost yet observed in biological channel encoding --- consistent with the large IB/IA ratio (\textasciitilde{}17) observed in AtHv1.

\textbf{Prediction.} Any other angiosperm Hv channel with IB/IA \textgreater{} 5 should encode K\_trap and T\_network in its resting state. Any channel with IB/IA $\approx$ 1 (gymnosperm, fungal, animal) should encode K\_mod and T\_bowtie. The IB/IA ratio is a direct experimental observable of the K primitive.

\textbf{Falsification.} An angiosperm Hv channel that shows IB/IA \textasciitilde{} 1 without a corresponding RSN would falsify the K\_trap/IB\_IA mapping. The IB/IA ratio for SmHv1 (\textasciitilde{}2.2, intermediate) should encode intermediate priming behavior --- either partial K\_trap or mixed RSN completeness.

\textbf{Status:} Tier I --- algebraically derived; d = 0.000 confirmed computationally in \texttt{hv1\_synthons.py}; chimera experimental data (Zhao et al. 2023 Nat. Commun.) directly supports encoding.

\begin{center}
\rule{0.5\textwidth}{0.4pt}
\end{center}

\subsection{P-56 $\cdot$ Hv1 Inhibitor Pathway Accessibility: K\_trap Blockade of Open-Channel Block}

\textbf{Primitive basis.} 2GBI is an open-channel blocker: it can only bind when Hv1 is in the depolarized/open state (Phi\_c, T\_bowtie). AtHv1 silent (K\_trap) never reaches the open state without mechanical priming --- the RSN prevents S4 translocation.

The critical distinction is not structural distance but \textbf{kinetic pathway accessibility}:

\begin{itemize}
    \item \texttt{tensor(2GBI, AtHv1\_silent)}: 2GBI is structurally close to AtHv1\_silent ($d = 1.0$) because both encode T\_network and P\_pm\_pseudo. But the binding pathway (channel opening) is blocked by K\_trap.
    \item \texttt{tensor(HIF, AtHv1\_primed)}: HIF ($d = 2.3$ from AtHv1\_primed) has a higher structural distance, but the pathway is now \emph{open} (K\_trap peeled).
\end{itemize}

\textbf{The primitive claim.} Structural tuple distance ($d$) measures \emph{encoding compatibility}, not \emph{pathway accessibility}. For open-channel blockers, the relevant quantity is:

\[\xi_{CP}^\text{effective} = \xi_{CP}^\text{binding} + \xi_{CP}^\text{opening pathway}\]

When K\_trap blocks the opening pathway, $\xi_{CP}^\text{opening} \to \infty$ --- the inhibitor cannot bind regardless of structural complementarity.

\textbf{Prediction.} Experimental test: measure IC50 of 2GBI against:

\begin{enumerate}
    \item AtHv1\_silent (no mechanical stimulus) $\rightarrow$ IC50 very high or unmeasurable (pathway blocked)
    \item AtHv1\_primed (post-mechanical stimulus) $\rightarrow$ IC50 measurable, comparable to PsHv1
\end{enumerate}

The ratio IC50(silent) / IC50(primed) for 2GBI should exceed the ratio for HIF, because HIF's fluorinated phenyl targets F182 which is accessible even in partial-open states. HIF should show \emph{relatively better} priming-dependence than 2GBI.

\textbf{Note on P-56 revision.} The original P-56 formulation used tuple distance as a proxy for binding compatibility. The correct formulation recognizes that K\_trap creates a pathway blockade invisible to $d$ but visible in \xi\_CP(effective). The revised prediction above is the falsifiable form.

\textbf{Falsification.} 2GBI showing significant inhibition of AtHv1\_silent (no mechanical priming required) would falsify the K\_trap pathway blockade claim --- implying 2GBI can access the binding site without full channel opening.

\textbf{Status:} Tier III --- algebraically derived; pathway blockade consistent with open-channel block mechanism (Hong et al. 2013); direct AtHv1 silent/primed IC50 comparison not yet published.

\begin{center}
\rule{0.5\textwidth}{0.4pt}
\end{center}

\subsection{P-57 $\cdot$ F150A Mutation Encodes $\Gamma$\_$\lor$(BROAD) Binding Site: Endogenous Guanidinium Promiscuity}

\textbf{Primitive basis.} In wild-type Hv1, the binding site around F150 encodes $\Gamma$\_$\land$(SELECTIVE): the condensed phenyl ring at F150 + F182 define tight geometric constraints that select for the specific 2GBI/ABI scaffold. The F150A mutation eliminates F150's steric contribution, creating a rearranged cavity where F182 moves closer to the vestibule.

This rearrangement changes the interaction grammar primitive of the binding site from $\Gamma$\_$\land$(SELECTIVE) to $\Gamma$\_$\lor$(BROAD): the cavity now has \emph{higher tolerance} for diverse chemical scaffolds (ABI outperforms 2GBI in F150A; ABIF₃ becomes potent where it was penalized in WT).

\textbf{Algebraic encoding.} The F150A binding site:

\[\langle D_\wedge; T_\in; R_\supset; P_{\pm\psi}; F_\eth; K_\text{fast}; G_\beth; \Gamma_\vee(\text{BROAD}); \Phi_\text{sub}; \Omega_0 \rangle\]

K\_fast (local rearrangement fast once the mutation is present); $\Gamma$\_$\lor$(BROAD) (permissive to structural diversity).

\textbf{Prediction.} A $\Gamma$\_$\lor$(BROAD) binding site is not selective for synthetic guanidinium compounds --- it should also accommodate structurally diverse endogenous guanidinium-containing molecules with comparable affinity. Candidates:

\begin{itemize}
    \item Agmatine (decarboxylated arginine, guanidinium group, endogenous in neurons)
    \item Arginine (the natural amino acid; guanidinium side chain)
    \item Creatine (guanidino-acetic acid, abundant in excitable cells)
    \item Homoarginine (lysine-derived guanidinium)
\end{itemize}

Any of these should show IC50 against Hv1 F150A in the nanomolar--micromolar range, while showing negligible inhibition of WT Hv1. The selectivity ratio IC50(WT) / IC50(F150A) should exceed 100-fold (matching the 2GBI shift from 38 $\mu$M to 118 nM).

\textbf{Broader implication.} If the F150 locus is naturally polymorphic in any cancer cell line or splice variant, those channels may show altered pharmacology toward endogenous guanidinium metabolites --- potentially linking cellular arginine metabolism to Hv1 activity in a new way.

\textbf{Falsification.} Agmatine or creatine showing IC50 \textgreater{} 1 mM against Hv1 F150A would falsify the $\Gamma$\_$\lor$(BROAD) interpretation, implying that the 2GBI-site selectivity is maintained despite the rearrangement.

\textbf{Status:} Tier III --- algebraically derived from F150A mutant data (Zhao et al. 2021 JGP); endogenous ligand screen not yet published.

\begin{center}
\rule{0.5\textwidth}{0.4pt}
\end{center}

\subsection{Summary Table}

\begin{table}[htbp]
\centering
\small
\rowcolors{1}{white}{white}
\begin{tabularx}{\textwidth}{>{\centering\arraybackslash}X >{\centering\arraybackslash}X >{\centering\arraybackslash}X >{\centering\arraybackslash}X >{\centering\arraybackslash}X}
\toprule
\rowcolor{gray!25}\textbf{ID} & \textbf{Prediction} & \textbf{Tier} & \textbf{Status} & \textbf{Key ref.} \\
\midrule
\rowcolors{1}{white}{gray!10}
P-1 & F-floor ratchet: 6/6 directions + asymmetric topology & I &  6/6 + ratchet irreversibility confirmed & Kim 2001; Assaf 2015 \\
P-2 & Soai $\rightarrow$ Frank bifurcation (Factor 7, score 0.920) & I &  confirmed & Soai 1995; Gridnev 2010 \\
P-3 & Proline-aldol ee = 70--85\% from F\_cycle & I &  74\% exp. & Blackmond 2004 \\
P-4 & Ice VI K\_fast $\rightarrow$ multiple ordered descendants & I &  confirmed (ice XV + XIX) & Yamane 2021; Salzmann 2009 \\
P-5 & H-bond dimer F ordering: AA \textgreater{} AA$\cdot$amide \textgreater{} amide & II &  DFT $\Delta$E ratio 1.9 & Transformation \#1 \\
P-6 & Triple H-bond induction superlinear 2.5--3.5$\times$ & II &  SAPT2+ confirmed & Transformation \#5 \\
P-7 & Halogen \textgreater{} chalcogen F: $\Delta$E ratio 1.91, V\_max 165/105 & II &  DFT + ESP confirmed & Transformation \#7 \\
P-8 & Chelate G\_beth $\rightarrow$ G\_aleph: +49 kJ/mol gain & II &  DFT confirmed & Transformation \#3 \\
P-9 & Imine K\_slow $\rightarrow$ K\_mod at aqueous barrier crossing & II &  DFT confirmed & Transformation \#4 \\
P-10 & Axiom 1 as classical boundary detector & II &  quantum particle series & V-5 \\
P-11 & $\Omega$ lattice semantics (tensor, meet, join) & II &  15 operations confirmed & v0.4.0 catalog \\
P-12 & K\_trap$\rightarrow$K\_MBL: +2.303 nat universal & II &  perturbation sweep & v0.4.0 quantum cluster \\
P-13 & Phase boundary d $\approx$ 9.52 (matter/particle) & II &  Ward + MDS confirmed & \texttt{syncon phase-diagram} \\
P-14 & Rotaxane steric cliff: K\_mod vs K\_trap sub-Å & II &  Groppi 2020 proxy & Transformation \#8 \\
P-15 & 8 algebra properties (F-floor, tensor bottleneck, \ldots{}) & II &  18/20 designs & V-6 \\
P-20 & $\lambda$ = primitive matching fraction; E{[}frac{]}=0.3023$\approx$$\lambda$\_fixed & II &  idempotency + catalog mean & Tensor algebra \\
P-21 & F-tier values = integer Boltzmann ratios (2:3, 3:1, 19:1) & II &  logit(3/4)=ln3, logit(19/20)=ln19 exact & Fidelity model \\
P-22 & $\Omega$ derivable from \{T,K,D,$\Gamma$,G\}; 5-rule, 0 mismatches & II &  32/32 synthons & Occam Target 3 \\
P-15b & K\_trap$\rightarrow$K\_MBL energy cost measurable (experimental target) & III & ⏳ $\alpha$-RuCl₃/InAs/FQH/SrCu₂O₃ & V-6 \\
P-23 & SM lift blocked (G=LOCAL); tensor(SM,QG) $\Phi$=Phi\_c; 4 CONFLICTS; no path & III & ⏳ encoding-conditional & §XVI \\
P-24 & G partitions gravity theories; tensor(GR,SM)$\rightarrow$F\_eth; AS $\Phi$\_c blocked; Hořava worst; Causal Set isolated; Verlinde$\rightarrow$T\_cup & II &  algebraically derived & §XVII \\
P-25 & \xi\_CP(BH)=0 at T\_H; mass-independent \eta\_CP; I\textasciitilde{}A from G\_aleph+T\_□□ & III & ⏳ analogue gravity systems & §XVIII.1 \\
P-26 & Factorization failure = frac=1 limit of P-20 tensor; \xi\_ens=max(\xi\_1,\xi\_2) & II/III &  algebraically; Type III extension pending & §XVIII.2 \\
P-27 & ER=EPR $\leftrightarrow$ R\_⇔\equiv$R_{\supseteq}$ at G\_aleph; extended Axiom 5; \xi\_CP(ER)=\xi\_CP(EPR) & III & ⏳ conditional on ER=EPR proof & §XVIII.3 \\
P-16 & Rotaxane $\Phi$\_c anchor (degeneracy\_strength $\geq$ 0.70) & III & ⏳ full scan pending & Transformation \#8 \\
P-17 & Ice XIX corr. length \textless{} ice XV (neutron diffraction) & III & ⏳ data not yet found & V-4 G\_beth vs G\_gimel \\
P-18 & Universality track d=2.303 in tuple-space MDS & III & ⏳ MBL synthons not yet added & Phase diagram \\
P-19 & \chi(T$\rightarrow$0)\textasciitilde{}T\^{}\{-$\gamma$\} for Factor-8 synthons, not TI & III & ⏳ not yet measured & Factor 8 \\
P-28 & $\beta$-sheet enzymes: higher active-site mutational tolerance ($d=2.80$ to scaffold) & II & ⏳ systematic alanine-scan comparison needed & Distance matrix, protein domain \\
P-29 & Complex assembly rate-limited by $K_\text{slow}$ subunit (tensor bottleneck) & II & ⏳ stopped-flow $k_\text{fold}$ vs. $k_\text{assemble}$ & Tensor bottleneck, V-6 \\
P-30 & Allosteric ON state: multi-timescale $R_{ex}$ spectrum (near-$\Phi_c$ broadening) & III & ⏳ CPMG relaxation dispersion on CAP & $\Phi_c$ assignment, Jacobian \\
P-31 & Rescue cost $<$ aggregation cost; $K$-targeting preferred over $F$-targeting & II & ⏳ head-to-head kinetics needed & Directed distance, F-floor \\
P-32 & Allosteric-interface chimera: binding OK, Hill $n_H \to 1$ & III & ⏳ chimera engineering + ITC & meet(allosteric, complex) conflicts \\
P-33 & Allosteric ($G_\gimel$) inhibitors more selective than orthosteric ($G_\beth$) & II & ⏳ selectivity panel comparison & $G$ primitive, KINOMEscan \\
P-34 & A$\beta$--$\alpha$-Syn cross-seeding $\approx$ homoseeding; both faster than tau & II &  consistent with Guo 2013; matched-control test pending & $d=0.00$ identity \\
P-35 & $K$-targeting dissolution rank: A$\beta$ $<$ tau $\approx$ $\alpha$-Syn vs. $F$-targeting & II & ⏳ chaperone vs. denaturant assay & Directed distance asymmetry \\
P-36 & Bivalent GNF-2: $T_\perp \to T_\in$ lowers $\xi_{CP}$, improves selectivity & III & ⏳ medicinal chemistry synthesis & GNF-2 gap analysis \\
P-37 & $G_\beth \to$ binary resistance; $G_\gimel \to$ incremental resistance & II &  $G_\beth$ confirmed (imatinib T315I); $G_\gimel$ prediction pending & $G$ primitive, clinical data \\
P-38 & Switching energy rank: $\Delta H_\text{SMP} < \Delta H_\text{LC} < \Delta H_\text{colloidal}$ & II & ⏳ $d$-rank 1.70 \textless{} 3.10 \textless{} 5.10 confirmed; DSC comparison pending & Programmability pair distance \\
P-39 & Gel rescue requires $\Gamma$- or K-targeting; thermal alone fails & II &  path blocked algebraically; consistent with Hsp70 biology & F-floor theorem, path algebra \\
P-40 & T-conflict pairs $\rightarrow$ first-order transitions (latent heat, hysteresis) & II &  LC N$\rightarrow$I and colloidal melting both first-order; prediction confirmed & T-conflict in meet, path topology \\
P-41 & Actin-DNA collective motion only with ATP ($\Phi_c$ ATP-conditional) & III & ⏳ composite derived; actin-DNA active motion not tested vs. ATP depletion & Tensor product, $\Phi_c$ propagation \\
P-42 & Maximally versatile PM generically near-critical (dynamic floor = $\Phi_c$) & II &  lattice meet confirmed; consistent with cortex/cytoplasm criticality & Programmability lattice §7 \\
P-43 & Condensates implement mesoscale allostery; $n_H \approx 0.60/0.50 \approx 1.2$ & III & ⏳ catalog analogy $d=2.50$; Hill coefficient prediction novel & Cross-domain analogy, Varma score \\
P-44 & Colloidal crystals with $d < 3.0$ from TI support boundary states & II &  topological colloids confirmed (Rechtsman 2016); $d$-threshold extends range & Catalog analogy \\
P-45 & Topological DNA origami shows polynomial (not exponential) strand-failure scaling & III & ⏳ topological origami exists; error scaling vs. topology class not tested & $\Omega$ primitive, catalog analogy \\
P-46 & Jacobian: F dominant design lever ($> 40\%$ d-reduction from F alone) & III & ⏳ computed algebraically; Tg vs. crosslink-density experiment pending & Primitive Jacobian §1 \\
P-47 & Actin-DNA composites structurally incoherent without ATP ($D$ conflict) & III & ⏳ consistent with actin biology; specific actin-DNA stability comparison pending & meet conflicts, $D$ primitive \\
P-48 & Condensate gelation and amyloid fibrillization are the same primitive event ($d=0.00$); K-targeting preferred therapeutic for both & II & ⏳ $d=0.00$ algebraically exact; K-targeting superiority over F-targeting experimentally testable (Jacobian hierarchy $K > F$ by $>1.5\times$) & Tuple distance §PM+Protein \\
P-49 & $\Phi_c$ is a categorical phase label, fully decoupled from ordinal $F$, $K$, $G$; $F$=LOW systems can be critical if $T$=NETWORK & II & ⏳ catalog proof: \texttt{asymptotic\_safety\_reuter\_fp}; consistent with LLPS; systematic ordinal-vs-criticality screen pending & Decomposition kernel, \texttt{phi\_c\_skeleton} \\
P-50 & Condensate$\rightarrow$amyloid conversion algebraically requires an external $F$=HIGH nucleation seed; unseeded condensates show indefinite lag phase & II & ⏳ F-CONFLICT algebraically exact; consistent with seeded aggregation literature; matched-control seed-free test pending & cofactor(amyloid, condensate) F-CONFLICT \\
P-51 & Quantization residual $B = \{K_\text{trap}, G_\aleph, T_\text{braid}, \Phi_c, \Omega_\text{NA}\}$ + D-CONFLICT encodes background-dependence problem & III & ⏳ algebraically derived; consistent with canonical quantisation obstacles & cofactor(QG, GR) per-primitive \\
P-52 & GNF-2 co-administered with a supramolecular network scaffold closes the $D$+$T$ allosteric gap; PROTAC/bispecific adaptor strategy & III & ⏳ algebraically derived; consistent with GNF-2 allosteric mechanism; no co-administration study published & Drug panel cofactor residual \\
P-53 & Three stability regimes (thermodynamic / phase-protected / kinetically frozen) are algebraically distinguishable by peel cost profile; LLPS cooperativity $>$ DNA origami at matched $T_m$ & II & ⏳ algebraically exact; consistent with LLPS phenomenology; head-to-head DSC comparison pending & Peel cost accounting, §decomposition \\
P-54 & AdS/CFT holography = GR $\otimes$ $\langle G_\aleph \otimes \Phi_c \rangle$; any system with global correlation + boundary criticality should exhibit holographic-like features & III & ⏳ algebraically derived; consistent with standard AdS/CFT; analogue gravity test not yet constructed & cofactor(AdS/CFT, GR), §§XVIII--XIX \\
P-55 & $d(\text{AtHv1\_primed},\ \text{PsHv1}) = 0.000$; mechanical priming collapses AtHv1 to exact gymnosperm constitutive state; IB/IA ratio is direct K\_trap observable & I &  $d=0.000$ computed in \texttt{hv1\_synthons.py}; chimera experiments (Zhao 2023) confirm RSN is the only distinction; IB/IA\textasciitilde{}17 vs \textasciitilde{}1.5 & Hv1 channel encoding; peel(K\_trap) \\
P-56 & K\_trap pathway blockade makes 2GBI ineffective against AtHv1\_silent regardless of structural complementarity; IC50(silent)/IC50(primed) \textgreater{} 100-fold for 2GBI; HIF shows relatively less priming-dependence & III & ⏳ IC50 ratio not yet measured; pathway blockade consistent with open-channel block mechanism & K\_trap pathway accessibility; \xi\_CP(effective) \\
P-57 & F150A binding site encodes $\Gamma_\vee(\text{BROAD})$; endogenous guanidinium compounds (agmatine, creatine, arginine) should inhibit Hv1 F150A at nanomolar--micromolar IC50 with \textgreater{}100-fold selectivity over WT & III & ⏳ endogenous ligand screen not yet published; $\Gamma$ change algebraically derived from F150A mutant data & F150A mutant encoding; $\Gamma$ primitive \\
\bottomrule
\end{tabularx}
\end{table}

\textbar{} P-ARCH-1 \textbar{} A holographic-native AI architecture trained to 10T parameters will exhibit grokking/generalization behavior equivalent to a Transformer at 143T \textbar{} III \textbar{} ⏳ No holographic-native architecture at scale exists yet; prediction is falsifiable when $D_\text{holo}$ bulk-boundary models reach 10T \textbar{} $D$ primitive governs relational density per parameter; $D_\text{holo} \Rightarrow G_\aleph$ built-in \textbar{}
\textbar{} P-ARCH-2 \textbar{} State-space models (Mamba-class, $D_\infty$) will hit $\Phi_c$ at an intermediate parameter count --- lower than 143T but higher than 10T, scaling as $\sim 143\text{T} / \ln(N)$ \textbar{} III \textbar{} ⏳ SSMs at scale emerging; $\Phi_c$ threshold not yet measured \textbar{} $D_\infty$ logarithmic memory scaling vs. $D_{\wedge\triangle}$ linear \textbar{}
\textbar{} P-ARCH-3 \textbar{} Pre-$\Phi_c$ hallucination spike in LLMs is a $K_\text{trap}$ signature --- error rates should follow the same asymmetric asymptote ($\xi_{CP}$ divergence) as other $K_\text{trap}$ systems approaching a phase boundary \textbar{} II \textbar{} ⏳ error-rate curves at scale are consistent with this shape; direct $\xi_{CP}$ fit not yet done \textbar{} $K_\text{trap}$ barrier structure, $\xi_{CP}$ divergence near $\Phi_c$ \textbar{}
\textbar{} P-ARCH-4 \textbar{} The 143.48T parameter threshold is architecture-specific ($D_{\wedge\triangle}$); it corresponds to the human synaptic relational density $86\text{T} \times 5/3$ and has no universal significance --- scaling beyond it in Transformer topology yields diminishing returns without a grammar shift ($\Gamma_\to \to \Gamma_\text{Quantum}$) \textbar{} II \textbar{} ⏳ scaling law flattening above 100T consistent with prediction; $\Gamma$ shift not directly measurable yet \textbar{} Human-synaptic proxy derivation; $G/D$ ratio misalignment above threshold \textbar{}

\begin{center}
\rule{0.5\textwidth}{0.4pt}
\end{center}

\subsection{Methodological note}

''Purely primitive-derived'' is interpreted strictly: a prediction counts as primitive-derived if \textbf{no domain equation is inserted between the primitive assignment and the predicted outcome}. What is permitted:

\begin{itemize}
    \item Ordinal comparisons ($F_{\hbar}$ \textgreater{} F\_eth is not a domain equation; it is the definition of the F lattice)
    \item Axiom firing rules (Factor 7, Axiom 1) applied to encoded tuples
    \item Algebraic operations (tensor, meet, join, lift, path) applied to the monad stack
    \item The K\_fast threshold ($\Delta$G‡ \textless{} 60 kJ/mol) --- this is the primitive definition, not a domain insertion
\end{itemize}

What disqualifies a prediction from this list: inserting a specific rate equation, reaction coordinate, or force-field term that is not encoded in the primitive tuple itself.

The distinction matters for falsifiability: if a prediction fails, the primitive encoding is wrong --- not a separate physics model.

\begin{center}
\rule{0.5\textwidth}{0.4pt}
\end{center}

\emph{Maintained in SynthOmnicon repository. Next review: after Transformation \#8 full scan and ice XIX correlation length data.}


\end{document}
